
% Default to the notebook output style

    


% Inherit from the specified cell style.




    
\documentclass[11pt, a4paper, french
]{book}

    
    
    \usepackage[T1]{fontenc}
    % Nicer default font (+ math font) than Computer Modern for most use cases
    \usepackage{mathpazo}
    \usepackage[francais]{babel} 

    % Basic figure setup, for now with no caption control since it's done
    % automatically by Pandoc (which extracts ![](path) syntax from Markdown).
    \usepackage{graphicx}
    % We will generate all images so they have a width \maxwidth. This means
    % that they will get their normal width if they fit onto the page, but
    % are scaled down if they would overflow the margins.
    \makeatletter
    \def\maxwidth{\ifdim\Gin@nat@width>\linewidth\linewidth
    \else\Gin@nat@width\fi}
    \makeatother
    \let\Oldincludegraphics\includegraphics
    % Set max figure width to be 80% of text width, for now hardcoded.
    \renewcommand{\includegraphics}[1]{\Oldincludegraphics[width=.8\maxwidth]{#1}}
    % Ensure that by default, figures have no caption (until we provide a
    % proper Figure object with a Caption API and a way to capture that
    % in the conversion process - todo).
    \usepackage{caption}
    \DeclareCaptionLabelFormat{nolabel}{}
    \captionsetup{labelformat=nolabel}

    \usepackage{adjustbox} % Used to constrain images to a maximum size 
    \usepackage{xcolor} % Allow colors to be defined
    \usepackage{enumerate} % Needed for markdown enumerations to work
    \usepackage{geometry} % Used to adjust the document margins
    \usepackage{amsmath} % Equations
    \usepackage{amssymb} % Equations
    \usepackage{textcomp} % defines textquotesingle
    % Hack from http://tex.stackexchange.com/a/47451/13684:
    \AtBeginDocument{%
        \def\PYZsq{\textquotesingle}% Upright quotes in Pygmentized code
    }
    \usepackage{upquote} % Upright quotes for verbatim code
    \usepackage{eurosym} % defines \euro
    \usepackage[mathletters]{ucs} % Extended unicode (utf-8) support
    \usepackage[utf8x]{inputenc} % Allow utf-8 characters in the tex document
    \usepackage{fancyvrb} % verbatim replacement that allows latex
    \usepackage{grffile} % extends the file name processing of package graphics 
                         % to support a larger range 
    % The hyperref package gives us a pdf with properly built
    % internal navigation ('pdf bookmarks' for the table of contents,
    % internal cross-reference links, web links for URLs, etc.)
    \usepackage{hyperref}
    \usepackage{longtable} % longtable support required by pandoc >1.10
    \usepackage{booktabs}  % table support for pandoc > 1.12.2
    \usepackage[inline]{enumitem} % IRkernel/repr support (it uses the enumerate* environment)
    \usepackage[normalem]{ulem} % ulem is needed to support strikethroughs (\sout)
                                % normalem makes italics be italics, not underlines
    

    
    
    % Colors for the hyperref package
    \definecolor{urlcolor}{rgb}{0,.145,.698}
    \definecolor{linkcolor}{rgb}{.71,0.21,0.01}
    \definecolor{citecolor}{rgb}{.12,.54,.11}

    % ANSI colors
    \definecolor{ansi-black}{HTML}{3E424D}
    \definecolor{ansi-black-intense}{HTML}{282C36}
    \definecolor{ansi-red}{HTML}{E75C58}
    \definecolor{ansi-red-intense}{HTML}{B22B31}
    \definecolor{ansi-green}{HTML}{00A250}
    \definecolor{ansi-green-intense}{HTML}{007427}
    \definecolor{ansi-yellow}{HTML}{DDB62B}
    \definecolor{ansi-yellow-intense}{HTML}{B27D12}
    \definecolor{ansi-blue}{HTML}{208FFB}
    \definecolor{ansi-blue-intense}{HTML}{0065CA}
    \definecolor{ansi-magenta}{HTML}{D160C4}
    \definecolor{ansi-magenta-intense}{HTML}{A03196}
    \definecolor{ansi-cyan}{HTML}{60C6C8}
    \definecolor{ansi-cyan-intense}{HTML}{258F8F}
    \definecolor{ansi-white}{HTML}{C5C1B4}
    \definecolor{ansi-white-intense}{HTML}{A1A6B2}

    % commands and environments needed by pandoc snippets
    % extracted from the output of `pandoc -s`
    \providecommand{\tightlist}{%
      \setlength{\itemsep}{0pt}\setlength{\parskip}{0pt}}
    \DefineVerbatimEnvironment{Highlighting}{Verbatim}{commandchars=\\\{\}}
    % Add ',fontsize=\small' for more characters per line
    \newenvironment{Shaded}{}{}
    \newcommand{\KeywordTok}[1]{\textcolor[rgb]{0.00,0.44,0.13}{\textbf{{#1}}}}
    \newcommand{\DataTypeTok}[1]{\textcolor[rgb]{0.56,0.13,0.00}{{#1}}}
    \newcommand{\DecValTok}[1]{\textcolor[rgb]{0.25,0.63,0.44}{{#1}}}
    \newcommand{\BaseNTok}[1]{\textcolor[rgb]{0.25,0.63,0.44}{{#1}}}
    \newcommand{\FloatTok}[1]{\textcolor[rgb]{0.25,0.63,0.44}{{#1}}}
    \newcommand{\CharTok}[1]{\textcolor[rgb]{0.25,0.44,0.63}{{#1}}}
    \newcommand{\StringTok}[1]{\textcolor[rgb]{0.25,0.44,0.63}{{#1}}}
    \newcommand{\CommentTok}[1]{\textcolor[rgb]{0.38,0.63,0.69}{\textit{{#1}}}}
    \newcommand{\OtherTok}[1]{\textcolor[rgb]{0.00,0.44,0.13}{{#1}}}
    \newcommand{\AlertTok}[1]{\textcolor[rgb]{1.00,0.00,0.00}{\textbf{{#1}}}}
    \newcommand{\FunctionTok}[1]{\textcolor[rgb]{0.02,0.16,0.49}{{#1}}}
    \newcommand{\RegionMarkerTok}[1]{{#1}}
    \newcommand{\ErrorTok}[1]{\textcolor[rgb]{1.00,0.00,0.00}{\textbf{{#1}}}}
    \newcommand{\NormalTok}[1]{{#1}}
    
    % Additional commands for more recent versions of Pandoc
    \newcommand{\ConstantTok}[1]{\textcolor[rgb]{0.53,0.00,0.00}{{#1}}}
    \newcommand{\SpecialCharTok}[1]{\textcolor[rgb]{0.25,0.44,0.63}{{#1}}}
    \newcommand{\VerbatimStringTok}[1]{\textcolor[rgb]{0.25,0.44,0.63}{{#1}}}
    \newcommand{\SpecialStringTok}[1]{\textcolor[rgb]{0.73,0.40,0.53}{{#1}}}
    \newcommand{\ImportTok}[1]{{#1}}
    \newcommand{\DocumentationTok}[1]{\textcolor[rgb]{0.73,0.13,0.13}{\textit{{#1}}}}
    \newcommand{\AnnotationTok}[1]{\textcolor[rgb]{0.38,0.63,0.69}{\textbf{\textit{{#1}}}}}
    \newcommand{\CommentVarTok}[1]{\textcolor[rgb]{0.38,0.63,0.69}{\textbf{\textit{{#1}}}}}
    \newcommand{\VariableTok}[1]{\textcolor[rgb]{0.10,0.09,0.49}{{#1}}}
    \newcommand{\ControlFlowTok}[1]{\textcolor[rgb]{0.00,0.44,0.13}{\textbf{{#1}}}}
    \newcommand{\OperatorTok}[1]{\textcolor[rgb]{0.40,0.40,0.40}{{#1}}}
    \newcommand{\BuiltInTok}[1]{{#1}}
    \newcommand{\ExtensionTok}[1]{{#1}}
    \newcommand{\PreprocessorTok}[1]{\textcolor[rgb]{0.74,0.48,0.00}{{#1}}}
    \newcommand{\AttributeTok}[1]{\textcolor[rgb]{0.49,0.56,0.16}{{#1}}}
    \newcommand{\InformationTok}[1]{\textcolor[rgb]{0.38,0.63,0.69}{\textbf{\textit{{#1}}}}}
    \newcommand{\WarningTok}[1]{\textcolor[rgb]{0.38,0.63,0.69}{\textbf{\textit{{#1}}}}}
    
    
    % Define a nice break command that doesn't care if a line doesn't already
    % exist.
    \def\br{\hspace*{\fill} \\* }
    % Math Jax compatability definitions
    \def\gt{>}
    \def\lt{<}
    % Document parameters
    \title{Python 3}
    
    
    

    % Pygments definitions
    
\makeatletter
\def\PY@reset{\let\PY@it=\relax \let\PY@bf=\relax%
    \let\PY@ul=\relax \let\PY@tc=\relax%
    \let\PY@bc=\relax \let\PY@ff=\relax}
\def\PY@tok#1{\csname PY@tok@#1\endcsname}
\def\PY@toks#1+{\ifx\relax#1\empty\else%
    \PY@tok{#1}\expandafter\PY@toks\fi}
\def\PY@do#1{\PY@bc{\PY@tc{\PY@ul{%
    \PY@it{\PY@bf{\PY@ff{#1}}}}}}}
\def\PY#1#2{\PY@reset\PY@toks#1+\relax+\PY@do{#2}}

\expandafter\def\csname PY@tok@w\endcsname{\def\PY@tc##1{\textcolor[rgb]{0.73,0.73,0.73}{##1}}}
\expandafter\def\csname PY@tok@c\endcsname{\let\PY@it=\textit\def\PY@tc##1{\textcolor[rgb]{0.25,0.50,0.50}{##1}}}
\expandafter\def\csname PY@tok@cp\endcsname{\def\PY@tc##1{\textcolor[rgb]{0.74,0.48,0.00}{##1}}}
\expandafter\def\csname PY@tok@k\endcsname{\let\PY@bf=\textbf\def\PY@tc##1{\textcolor[rgb]{0.00,0.50,0.00}{##1}}}
\expandafter\def\csname PY@tok@kp\endcsname{\def\PY@tc##1{\textcolor[rgb]{0.00,0.50,0.00}{##1}}}
\expandafter\def\csname PY@tok@kt\endcsname{\def\PY@tc##1{\textcolor[rgb]{0.69,0.00,0.25}{##1}}}
\expandafter\def\csname PY@tok@o\endcsname{\def\PY@tc##1{\textcolor[rgb]{0.40,0.40,0.40}{##1}}}
\expandafter\def\csname PY@tok@ow\endcsname{\let\PY@bf=\textbf\def\PY@tc##1{\textcolor[rgb]{0.67,0.13,1.00}{##1}}}
\expandafter\def\csname PY@tok@nb\endcsname{\def\PY@tc##1{\textcolor[rgb]{0.00,0.50,0.00}{##1}}}
\expandafter\def\csname PY@tok@nf\endcsname{\def\PY@tc##1{\textcolor[rgb]{0.00,0.00,1.00}{##1}}}
\expandafter\def\csname PY@tok@nc\endcsname{\let\PY@bf=\textbf\def\PY@tc##1{\textcolor[rgb]{0.00,0.00,1.00}{##1}}}
\expandafter\def\csname PY@tok@nn\endcsname{\let\PY@bf=\textbf\def\PY@tc##1{\textcolor[rgb]{0.00,0.00,1.00}{##1}}}
\expandafter\def\csname PY@tok@ne\endcsname{\let\PY@bf=\textbf\def\PY@tc##1{\textcolor[rgb]{0.82,0.25,0.23}{##1}}}
\expandafter\def\csname PY@tok@nv\endcsname{\def\PY@tc##1{\textcolor[rgb]{0.10,0.09,0.49}{##1}}}
\expandafter\def\csname PY@tok@no\endcsname{\def\PY@tc##1{\textcolor[rgb]{0.53,0.00,0.00}{##1}}}
\expandafter\def\csname PY@tok@nl\endcsname{\def\PY@tc##1{\textcolor[rgb]{0.63,0.63,0.00}{##1}}}
\expandafter\def\csname PY@tok@ni\endcsname{\let\PY@bf=\textbf\def\PY@tc##1{\textcolor[rgb]{0.60,0.60,0.60}{##1}}}
\expandafter\def\csname PY@tok@na\endcsname{\def\PY@tc##1{\textcolor[rgb]{0.49,0.56,0.16}{##1}}}
\expandafter\def\csname PY@tok@nt\endcsname{\let\PY@bf=\textbf\def\PY@tc##1{\textcolor[rgb]{0.00,0.50,0.00}{##1}}}
\expandafter\def\csname PY@tok@nd\endcsname{\def\PY@tc##1{\textcolor[rgb]{0.67,0.13,1.00}{##1}}}
\expandafter\def\csname PY@tok@s\endcsname{\def\PY@tc##1{\textcolor[rgb]{0.73,0.13,0.13}{##1}}}
\expandafter\def\csname PY@tok@sd\endcsname{\let\PY@it=\textit\def\PY@tc##1{\textcolor[rgb]{0.73,0.13,0.13}{##1}}}
\expandafter\def\csname PY@tok@si\endcsname{\let\PY@bf=\textbf\def\PY@tc##1{\textcolor[rgb]{0.73,0.40,0.53}{##1}}}
\expandafter\def\csname PY@tok@se\endcsname{\let\PY@bf=\textbf\def\PY@tc##1{\textcolor[rgb]{0.73,0.40,0.13}{##1}}}
\expandafter\def\csname PY@tok@sr\endcsname{\def\PY@tc##1{\textcolor[rgb]{0.73,0.40,0.53}{##1}}}
\expandafter\def\csname PY@tok@ss\endcsname{\def\PY@tc##1{\textcolor[rgb]{0.10,0.09,0.49}{##1}}}
\expandafter\def\csname PY@tok@sx\endcsname{\def\PY@tc##1{\textcolor[rgb]{0.00,0.50,0.00}{##1}}}
\expandafter\def\csname PY@tok@m\endcsname{\def\PY@tc##1{\textcolor[rgb]{0.40,0.40,0.40}{##1}}}
\expandafter\def\csname PY@tok@gh\endcsname{\let\PY@bf=\textbf\def\PY@tc##1{\textcolor[rgb]{0.00,0.00,0.50}{##1}}}
\expandafter\def\csname PY@tok@gu\endcsname{\let\PY@bf=\textbf\def\PY@tc##1{\textcolor[rgb]{0.50,0.00,0.50}{##1}}}
\expandafter\def\csname PY@tok@gd\endcsname{\def\PY@tc##1{\textcolor[rgb]{0.63,0.00,0.00}{##1}}}
\expandafter\def\csname PY@tok@gi\endcsname{\def\PY@tc##1{\textcolor[rgb]{0.00,0.63,0.00}{##1}}}
\expandafter\def\csname PY@tok@gr\endcsname{\def\PY@tc##1{\textcolor[rgb]{1.00,0.00,0.00}{##1}}}
\expandafter\def\csname PY@tok@ge\endcsname{\let\PY@it=\textit}
\expandafter\def\csname PY@tok@gs\endcsname{\let\PY@bf=\textbf}
\expandafter\def\csname PY@tok@gp\endcsname{\let\PY@bf=\textbf\def\PY@tc##1{\textcolor[rgb]{0.00,0.00,0.50}{##1}}}
\expandafter\def\csname PY@tok@go\endcsname{\def\PY@tc##1{\textcolor[rgb]{0.53,0.53,0.53}{##1}}}
\expandafter\def\csname PY@tok@gt\endcsname{\def\PY@tc##1{\textcolor[rgb]{0.00,0.27,0.87}{##1}}}
\expandafter\def\csname PY@tok@err\endcsname{\def\PY@bc##1{\setlength{\fboxsep}{0pt}\fcolorbox[rgb]{1.00,0.00,0.00}{1,1,1}{\strut ##1}}}
\expandafter\def\csname PY@tok@kc\endcsname{\let\PY@bf=\textbf\def\PY@tc##1{\textcolor[rgb]{0.00,0.50,0.00}{##1}}}
\expandafter\def\csname PY@tok@kd\endcsname{\let\PY@bf=\textbf\def\PY@tc##1{\textcolor[rgb]{0.00,0.50,0.00}{##1}}}
\expandafter\def\csname PY@tok@kn\endcsname{\let\PY@bf=\textbf\def\PY@tc##1{\textcolor[rgb]{0.00,0.50,0.00}{##1}}}
\expandafter\def\csname PY@tok@kr\endcsname{\let\PY@bf=\textbf\def\PY@tc##1{\textcolor[rgb]{0.00,0.50,0.00}{##1}}}
\expandafter\def\csname PY@tok@bp\endcsname{\def\PY@tc##1{\textcolor[rgb]{0.00,0.50,0.00}{##1}}}
\expandafter\def\csname PY@tok@fm\endcsname{\def\PY@tc##1{\textcolor[rgb]{0.00,0.00,1.00}{##1}}}
\expandafter\def\csname PY@tok@vc\endcsname{\def\PY@tc##1{\textcolor[rgb]{0.10,0.09,0.49}{##1}}}
\expandafter\def\csname PY@tok@vg\endcsname{\def\PY@tc##1{\textcolor[rgb]{0.10,0.09,0.49}{##1}}}
\expandafter\def\csname PY@tok@vi\endcsname{\def\PY@tc##1{\textcolor[rgb]{0.10,0.09,0.49}{##1}}}
\expandafter\def\csname PY@tok@vm\endcsname{\def\PY@tc##1{\textcolor[rgb]{0.10,0.09,0.49}{##1}}}
\expandafter\def\csname PY@tok@sa\endcsname{\def\PY@tc##1{\textcolor[rgb]{0.73,0.13,0.13}{##1}}}
\expandafter\def\csname PY@tok@sb\endcsname{\def\PY@tc##1{\textcolor[rgb]{0.73,0.13,0.13}{##1}}}
\expandafter\def\csname PY@tok@sc\endcsname{\def\PY@tc##1{\textcolor[rgb]{0.73,0.13,0.13}{##1}}}
\expandafter\def\csname PY@tok@dl\endcsname{\def\PY@tc##1{\textcolor[rgb]{0.73,0.13,0.13}{##1}}}
\expandafter\def\csname PY@tok@s2\endcsname{\def\PY@tc##1{\textcolor[rgb]{0.73,0.13,0.13}{##1}}}
\expandafter\def\csname PY@tok@sh\endcsname{\def\PY@tc##1{\textcolor[rgb]{0.73,0.13,0.13}{##1}}}
\expandafter\def\csname PY@tok@s1\endcsname{\def\PY@tc##1{\textcolor[rgb]{0.73,0.13,0.13}{##1}}}
\expandafter\def\csname PY@tok@mb\endcsname{\def\PY@tc##1{\textcolor[rgb]{0.40,0.40,0.40}{##1}}}
\expandafter\def\csname PY@tok@mf\endcsname{\def\PY@tc##1{\textcolor[rgb]{0.40,0.40,0.40}{##1}}}
\expandafter\def\csname PY@tok@mh\endcsname{\def\PY@tc##1{\textcolor[rgb]{0.40,0.40,0.40}{##1}}}
\expandafter\def\csname PY@tok@mi\endcsname{\def\PY@tc##1{\textcolor[rgb]{0.40,0.40,0.40}{##1}}}
\expandafter\def\csname PY@tok@il\endcsname{\def\PY@tc##1{\textcolor[rgb]{0.40,0.40,0.40}{##1}}}
\expandafter\def\csname PY@tok@mo\endcsname{\def\PY@tc##1{\textcolor[rgb]{0.40,0.40,0.40}{##1}}}
\expandafter\def\csname PY@tok@ch\endcsname{\let\PY@it=\textit\def\PY@tc##1{\textcolor[rgb]{0.25,0.50,0.50}{##1}}}
\expandafter\def\csname PY@tok@cm\endcsname{\let\PY@it=\textit\def\PY@tc##1{\textcolor[rgb]{0.25,0.50,0.50}{##1}}}
\expandafter\def\csname PY@tok@cpf\endcsname{\let\PY@it=\textit\def\PY@tc##1{\textcolor[rgb]{0.25,0.50,0.50}{##1}}}
\expandafter\def\csname PY@tok@c1\endcsname{\let\PY@it=\textit\def\PY@tc##1{\textcolor[rgb]{0.25,0.50,0.50}{##1}}}
\expandafter\def\csname PY@tok@cs\endcsname{\let\PY@it=\textit\def\PY@tc##1{\textcolor[rgb]{0.25,0.50,0.50}{##1}}}

\def\PYZbs{\char`\\}
\def\PYZus{\char`\_}
\def\PYZob{\char`\{}
\def\PYZcb{\char`\}}
\def\PYZca{\char`\^}
\def\PYZam{\char`\&}
\def\PYZlt{\char`\<}
\def\PYZgt{\char`\>}
\def\PYZsh{\char`\#}
\def\PYZpc{\char`\%}
\def\PYZdl{\char`\$}
\def\PYZhy{\char`\-}
\def\PYZsq{\char`\'}
\def\PYZdq{\char`\"}
\def\PYZti{\char`\~}
% for compatibility with earlier versions
\def\PYZat{@}
\def\PYZlb{[}
\def\PYZrb{]}
\makeatother


    % Exact colors from NB
    \definecolor{incolor}{rgb}{0.0, 0.0, 0.5}
    \definecolor{outcolor}{rgb}{0.545, 0.0, 0.0}



    
    % Prevent overflowing lines due to hard-to-break entities
    \sloppy 
    % Setup hyperref package
    \hypersetup{
      breaklinks=true,  % so long urls are correctly broken across lines
      colorlinks=true,
      urlcolor=urlcolor,
      linkcolor=linkcolor,
      citecolor=citecolor,
      }
    % Slightly bigger margins than the latex defaults
    
    \geometry{verbose,tmargin=1in,bmargin=1in,lmargin=1in,rmargin=1in}
    
    

    \begin{document}
    
    
    \maketitle
    
    \chapter{Semaine 1}
        
    
    
    

    

    \hypertarget{installer-la-distribution-standard-python}{%
\section{Installer la distribution standard
Python}\label{installer-la-distribution-standard-python}}

    \hypertarget{compluxe9ment---niveau-basique}{%
\subsection{Complément - niveau
basique}\label{compluxe9ment---niveau-basique}}

    Ce complément a pour but de vous donner quelques guides pour
l'installation de la distribution standard Python 3.

Notez bien qu'il ne s'agit ici que d'indications, il existe de
nombreuses façons de procéder.

En cas de souci, commencez par chercher par vous-même, sur Google ou
autre, une solution à votre problème~; pensez également à utiliser le
forum du cours.

    Le point important est de \textbf{bien vérifier le numéro de version} de
votre installation qui doit être \textbf{au moins 3.6}

    \hypertarget{sachez-uxe0-qui-vous-parlez}{%
\subsection{Sachez à qui vous
parlez}\label{sachez-uxe0-qui-vous-parlez}}

    Mais avant qu'on n'avance sur l'installation proprement dite, il nous
faut insister sur un point qui déroute parfois les débutants. On a
parfois besoin de recourir à l'emploi d'un terminal, surtout justement
pendant la phase d'installation.

Lorsque c'est le cas, il est important de bien distinguer~:

\begin{itemize}
\tightlist
\item
  les cas où on s'adresse \textbf{au terminal} (en jargon, on dit le
  \emph{shell}),
\item
  et les cas où on s'adresse à \textbf{l'interpréteur Python}.
\end{itemize}

C'est très important car ces deux programmes ne parlent \textbf{pas} du
tout le \textbf{même langage} ! Il peut arriver au début qu'on écrive
une commande juste, mais au mauvais interlocuteur, et cela peut être
source de frustration. Essayons de bien comprendre ce point.

    \hypertarget{le-terminal}{%
\subsubsection{Le terminal}\label{le-terminal}}

Je peux dire que je parle à mon \textbf{terminal} quand l'invite de
commande (en jargon on dit le \emph{prompt}) \textbf{se termine par un
dollar \texttt{\$}} - ou un simple chevron \texttt{\textgreater{}} sur
Windows

Par exemple~sur un mac :

\begin{Shaded}
\begin{Highlighting}[frame=lines,framerule=0.6mm,rulecolor=\color{asisframecolor}]
\ExtensionTok{~/git/flotpython/w1}\NormalTok{ $ }
\end{Highlighting}
\end{Shaded}

Ou sur Windows~:

\begin{Shaded}
\begin{Highlighting}[frame=lines,framerule=0.6mm,rulecolor=\color{asisframecolor}]
\ExtensionTok{C}\NormalTok{:\textbackslash{}Users}\OperatorTok{>}
\end{Highlighting}
\end{Shaded}

    \hypertarget{linterpruxe8te-python}{%
\subsubsection{L'interprète Python}\label{linterpruxe8te-python}}

À partir du terminal, je peux lancer un \textbf{interpréteur Python},
qui se reconnaît car son prompt est fait de \textbf{3 chevrons
\texttt{\textgreater{}\textgreater{}\textgreater{}}}

\begin{Shaded}
\begin{Highlighting}[frame=lines,framerule=0.6mm,rulecolor=\color{asisframecolor}]
\ExtensionTok{~/git/flotpython/w1}\NormalTok{ $ python3}
\ExtensionTok{Python}\NormalTok{ 3.7.0 (default, Jun 29 2018, 20:14:27)}
\NormalTok{[}\ExtensionTok{Clang}\NormalTok{ 9.0.0 (clang-900.0.39.2)] }\ExtensionTok{on}\NormalTok{ darwin}
\ExtensionTok{Type} \StringTok{"help"}\NormalTok{, }\StringTok{"copyright"}\NormalTok{, }\StringTok{"credits"}\NormalTok{ or }\StringTok{"license"}\NormalTok{ for more information.}
\OperatorTok{>>>}
\end{Highlighting}
\end{Shaded}

    Pour sortir de l'interpréteur Python, et retourner au terminal,
j'utilise la fonction Python \texttt{exit()}~:

\begin{Shaded}
\begin{Highlighting}[frame=lines,framerule=0.6mm,rulecolor=\color{asisframecolor}]
\ExtensionTok{~/git/flotpython/w1}\NormalTok{ $ python3}
\OperatorTok{>>>} \ExtensionTok{20}\NormalTok{ * 60}
\ExtensionTok{1200}
\OperatorTok{>>>} \FunctionTok{exit()}
\ExtensionTok{~/git/flotpython/w1}\NormalTok{ $ python3}
\end{Highlighting}
\end{Shaded}

    \hypertarget{les-erreurs-typiques}{%
\subsubsection{Les erreurs typiques}\label{les-erreurs-typiques}}

Gardez bien cette distinction présente à l'esprit, lorsque vous lisez la
suite. Voici quelques symptômes habituels de ce qu'on obtient si on se
trompe.

Par exemple, la commande \texttt{python3\ -V} est une commande qui
s'adresse au terminal~; c'est pourquoi nous la faisons précéder d'un
dollar \texttt{\$}.

Si vous essayez de la taper alors que vous êtes déjà dans un
interpréteur python - ou sous IDLE d'ailleurs -, vous obtenez un message
d'erreur de ce genre :

\begin{Shaded}
\begin{Highlighting}[frame=lines,framerule=0.6mm,rulecolor=\color{asisframecolor}]
\OperatorTok{>>>} \ExtensionTok{python3}\NormalTok{ -V}
\ExtensionTok{Traceback}\NormalTok{ (most recent call last)}\BuiltInTok{:}
  \ExtensionTok{File} \StringTok{"<stdin>"}\NormalTok{, line 1, in }\OperatorTok{<}\NormalTok{module}\OperatorTok{>}
\ExtensionTok{NameError}\NormalTok{: name }\StringTok{'python3'}\NormalTok{ is not defined}
\end{Highlighting}
\end{Shaded}

    Réciproquement, si vous essayez de taper du Python directement dans un
terminal, ça se passe mal aussi, forcément. Par exemple sur Mac, avec
des fragments Python tout simples~:

\begin{Shaded}
\begin{Highlighting}[frame=lines,framerule=0.6mm,rulecolor=\color{asisframecolor}]
\ExtensionTok{~/git/flotpython/w1}\NormalTok{ $ import math}
\ExtensionTok{-bash}\NormalTok{: import: command not found}
\ExtensionTok{~/git/flotpython/w1}\NormalTok{ $ 30 * 60}
\ExtensionTok{-bash}\NormalTok{: 30: command not found}
\ExtensionTok{~/git/flotpython/w1}\NormalTok{ $ foo = 30 * 60}
\ExtensionTok{-bash}\NormalTok{: foo: command not found}
\end{Highlighting}
\end{Shaded}

    \hypertarget{digression---coexistence-de-python2-et-python3}{%
\subsection{Digression - coexistence de Python2 et
Python3}\label{digression---coexistence-de-python2-et-python3}}

    Avant l'arrivée de la version 3 de Python, les choses étaient simples,
on exécutait un programme Python avec une seule commande
\texttt{python}. Depuis 2014-2015, maintenant que les deux versions de
Python coexistent, il est nécessaire d'adopter une convention qui
permette d'installer les deux langages sous des noms qui sont
non-ambigus.

C'est pourquoi actuellement, on trouve \textbf{le plus souvent} la
convention suivante sous Linux et macOS~:

\begin{itemize}
\item
  \texttt{python3} est pour exécuter les programmes en Python-3~; du
  coup on trouve alors également les commandes comme \texttt{idle3} pour
  lancer IDLE, et par exemple \texttt{pip3} pour le gestionnaire de
  paquets (voir ci-dessous)~;
\item
  \texttt{python2} est pour exécuter les programmes en Python-2, avec
  typiquement \texttt{idle2} et \texttt{pip2}~;
\item
  enfin selon les systèmes, la commande \texttt{python} tout court est
  un alias pour \texttt{python2} ou \texttt{python3}. De plus en plus
  souvent, par défaut \texttt{python} désigne \texttt{python3}.
\end{itemize}

à titre d'illustration, voici ce que j'obtiens sur mon mac~:

\begin{Shaded}
\begin{Highlighting}[frame=lines,framerule=0.6mm,rulecolor=\color{asisframecolor}]
\NormalTok{$ }\ExtensionTok{python3}\NormalTok{ -V}
\ExtensionTok{Python}\NormalTok{ 3.6.2}
\NormalTok{$ }\ExtensionTok{python2}\NormalTok{ -V}
\ExtensionTok{Python}\NormalTok{ 2.7.13}
\NormalTok{$ }\ExtensionTok{python}\NormalTok{ -V}
\ExtensionTok{Python}\NormalTok{ 3.6.2}
\end{Highlighting}
\end{Shaded}

Sous Windows, vous avez un lanceur qui s'appelle \texttt{py}. Par
défaut, il lance la version de Python la plus récente installée, mais
vous pouvez spécifier une version spécifique de la manière suivante~:

\begin{Shaded}
\begin{Highlighting}[frame=lines,framerule=0.6mm,rulecolor=\color{asisframecolor}]
\ExtensionTok{C}\NormalTok{:}\DataTypeTok{\textbackslash{}>}\NormalTok{ py -2.7}
\end{Highlighting}
\end{Shaded}

pour lancer, par exemple, Python en version 2.7. Vous trouverez toute la
documentation nécessaire pour Windows sur cette page (en anglais)~:
https://docs.python.org/3/using/windows.html

Pour éviter d'éventuelles confusions, nous précisons toujours
\texttt{python3} dans le cours.

    \hypertarget{installation-de-base}{%
\subsection{Installation de base}\label{installation-de-base}}

    \hypertarget{vous-utilisez-windows}{%
\subsubsection{Vous utilisez Windows}\label{vous-utilisez-windows}}

    La méthode recommandée sur Windows est de partir de la page
https://www.python.org/download où vous trouverez un programme
d'installation qui contient tout ce dont vous aurez besoin pour suivre
le cours.

Pour vérifier que vous êtes prêt, il vous faut lancer IDLE (quelque part
dans le menu Démarrer) et vérifier le numéro de version.

    \hypertarget{vous-utilisez-macos}{%
\subsubsection{Vous utilisez macOS}\label{vous-utilisez-macos}}

    Ici encore, la méthode recommandée est de partir de la page
https://www.python.org/download et d'utiliser le programme
d'installation.

Sachez aussi, si vous utilisez déjà MacPorts (https://www.macports.org),
que vous pouvez également utiliser cet outil pour installer, par exemple
Python 3.6, avec la commande

    \begin{Shaded}
\begin{Highlighting}[frame=lines,framerule=0.6mm,rulecolor=\color{asisframecolor}]
\NormalTok{$ }\FunctionTok{sudo}\NormalTok{ port install python36}
\end{Highlighting}
\end{Shaded}

    \hypertarget{vous-utilisez-linux}{%
\subsubsection{Vous utilisez Linux}\label{vous-utilisez-linux}}

    Dans ce cas il est très probable que Python-3.x soit déjà disponible sur
votre machine. Pour vous en assurer, essayez de lancer la commande
\texttt{python3} dans un terminal.

    \hypertarget{rhel-fedora}{%
\subparagraph{RHEL / Fedora}\label{rhel-fedora}}

    Voici par exemple ce qu'on obtient depuis un terminal sur une machine
installée en Fedora-20~:

    \begin{Shaded}
\begin{Highlighting}[frame=lines,framerule=0.6mm,rulecolor=\color{asisframecolor}]
\NormalTok{$ }\ExtensionTok{python3}
\ExtensionTok{Python}\NormalTok{ 3.6.2 (default, Jul 20 2017, 12:30:02)}
\NormalTok{[}\ExtensionTok{GCC}\NormalTok{ 6.3.1 20161221 (Red Hat 6.3.1-1)] }\ExtensionTok{on}\NormalTok{ linux}
\ExtensionTok{Type} \StringTok{"help"}\NormalTok{, }\StringTok{"copyright"}\NormalTok{, }\StringTok{"credits"}\NormalTok{ or }\StringTok{"license"}\NormalTok{ for more information.}
\OperatorTok{>>>} \FunctionTok{exit()}
\end{Highlighting}
\end{Shaded}

    \textbf{Vérifiez bien le numéro de version} qui doit être en 3.\emph{x}.
Si vous obtenez un message du style
\texttt{python3:\ command\ not\ found} utilisez \texttt{dnf}
(anciennement connu sous le nom de \texttt{yum}) pour installer le rpm
\texttt{python3} comme ceci~:

    \begin{Shaded}
\begin{Highlighting}[frame=lines,framerule=0.6mm,rulecolor=\color{asisframecolor}]
\NormalTok{$ }\FunctionTok{sudo}\NormalTok{ dnf install python3}
\end{Highlighting}
\end{Shaded}

    S'agissant d'\texttt{idle}, l'éditeur que nous utilisons dans le cours
(optionnel si vous êtes familier avec un éditeur de texte), vérifiez sa
présence comme ceci~:

    \begin{Shaded}
\begin{Highlighting}[frame=lines,framerule=0.6mm,rulecolor=\color{asisframecolor}]
\NormalTok{$ }\BuiltInTok{type}\NormalTok{ idle3}
\ExtensionTok{idle}\NormalTok{ is hashed (/usr/bin/idle3)}
\end{Highlighting}
\end{Shaded}

    Ici encore, si la commande n'est pas disponible vous pouvez l'installer
avec~:

    \begin{Shaded}
\begin{Highlighting}[frame=lines,framerule=0.6mm,rulecolor=\color{asisframecolor}]
\NormalTok{$ }\FunctionTok{sudo}\NormalTok{ yum install python3-tools}
\end{Highlighting}
\end{Shaded}

    \hypertarget{debian-ubuntu}{%
\subparagraph{Debian / Ubuntu}\label{debian-ubuntu}}

    Ici encore, Python-2.7 est sans doute déjà disponible. Procédez comme
ci-dessus, voici un exemple recueilli dans un terminal sur une machine
installée en Ubuntu-14.04/trusty~:

    \begin{Shaded}
\begin{Highlighting}[frame=lines,framerule=0.6mm,rulecolor=\color{asisframecolor}]
\NormalTok{$ }\ExtensionTok{python3}
\ExtensionTok{Python}\NormalTok{ 3.6.2 (default, Jul 20 2017, 12:30:02)}
\NormalTok{[}\ExtensionTok{GCC}\NormalTok{ 6.3.1 20161221 (Red Hat 6.3.1-1)] }\ExtensionTok{on}\NormalTok{ linux}
\ExtensionTok{Type} \StringTok{"help"}\NormalTok{, }\StringTok{"copyright"}\NormalTok{, }\StringTok{"credits"}\NormalTok{ or }\StringTok{"license"}\NormalTok{ for more information.}
\OperatorTok{>>>} \FunctionTok{exit()}
\end{Highlighting}
\end{Shaded}

    Pour installer Python~:

    \begin{Shaded}
\begin{Highlighting}[frame=lines,framerule=0.6mm,rulecolor=\color{asisframecolor}]
\NormalTok{$ }\FunctionTok{sudo}\NormalTok{ apt-get install python3}
\end{Highlighting}
\end{Shaded}

    Pour installer idle~:

    \begin{Shaded}
\begin{Highlighting}[frame=lines,framerule=0.6mm,rulecolor=\color{asisframecolor}]
\NormalTok{$ }\FunctionTok{sudo}\NormalTok{ apt-get install idle3}
\end{Highlighting}
\end{Shaded}

    \hypertarget{installation-de-bibliothuxe8ques-compluxe9mentaires}{%
\subsubsection{Installation de bibliothèques
complémentaires}\label{installation-de-bibliothuxe8ques-compluxe9mentaires}}

    Il existe un outil très pratique pour installer des bibliothèques
Python, il s'appelle \texttt{pip3}, qui est documenté ici
https://pypi.python.org/pypi/pip

    Sachez aussi, si par ailleurs vous utilisez un gestionnaire de paquets
comme \texttt{rpm} sur RHEL, \texttt{apt-get} sur Debian, ou
\texttt{port} sur macOS, que de nombreux paquets sont également
disponibles au travers de ces outils.

    \hypertarget{anaconda}{%
\subsubsection{Anaconda}\label{anaconda}}

    Sachez qu'il existe beaucoup de distributions alternatives qui incluent
Python~; parmi elles, la plus populaire est sans aucun doute
\href{https://www.anaconda.com/}{Anaconda}, qui contient un grand nombre
de bibliothèques de calcul scientifique, et également d'ailleurs Jupyter
pour travailler nativement sur des notebooks au format \texttt{.ipynb}.

Anaconda vient avec son propre gestionnaire de paquets pour
l'installation de bibliothèques supplémentaires qui s'appelle
\texttt{conda}.


    % Add a bibliography block to the postdoc
    
    
    

        
    
    
    

    

    \hypertarget{un-peu-de-lecture}{%
\section{Un peu de lecture}\label{un-peu-de-lecture}}

    \hypertarget{compluxe9ment---niveau-basique}{%
\subsection{Complément - niveau
basique}\label{compluxe9ment---niveau-basique}}

    \hypertarget{mise-uxe0-jour-de-juillet-2018}{%
\subsubsection{Mise à jour de Juillet
2018}\label{mise-uxe0-jour-de-juillet-2018}}

    Le 12 Juillet 2018, Guido van Rossum
\href{https://lwn.net/Articles/759654/}{a annoncé qu'il quittait la
fonction de BDFL} qu'il occupait depuis près de trois décennies. Il
n'est pas tout à fait clair à ce stade comment va évoluer la gouvernance
de Python.

    \hypertarget{le-zen-de-python}{%
\subsubsection{Le Zen de Python}\label{le-zen-de-python}}

    Vous pouvez lire le ``Zen de Python'', qui résume la philosophie du
langage, en important le module \texttt{this} avec ce code~: (pour
exécuter ce code, cliquez dans la cellule de code, et faites au clavier
``Majuscule/Entrée'' ou ``Shift/Enter'')

    \begin{Verbatim}[commandchars=\\\{\},frame=single,framerule=0.3mm,rulecolor=\color{cellframecolor}]
{\color{incolor}In [{\color{incolor}1}]:} \PY{c+c1}{\PYZsh{} le Zen de Python}
        \PY{k+kn}{import} \PY{n+nn}{this}
\end{Verbatim}


    \begin{Verbatim}[commandchars=\\\{\},frame=single,framerule=0.3mm,rulecolor=\color{cellframecolor}]
The Zen of Python, by Tim Peters

Beautiful is better than ugly.
Explicit is better than implicit.
Simple is better than complex.
Complex is better than complicated.
Flat is better than nested.
Sparse is better than dense.
Readability counts.
Special cases aren't special enough to break the rules.
Although practicality beats purity.
Errors should never pass silently.
Unless explicitly silenced.
In the face of ambiguity, refuse the temptation to guess.
There should be one-- and preferably only one --obvious way to do it.
Although that way may not be obvious at first unless you're Dutch.
Now is better than never.
Although never is often better than *right* now.
If the implementation is hard to explain, it's a bad idea.
If the implementation is easy to explain, it may be a good idea.
Namespaces are one honking great idea -- let's do more of those!
\end{Verbatim}

    \hypertarget{documentation}{%
\subsubsection{Documentation}\label{documentation}}

    \begin{itemize}
\item
  On peut commencer par citer
  l'\href{http://fr.wikipedia.org/wiki/Python_\%28langage\%29}{article
  de Wikipédia sur Python en français}.
\item
  La \href{https://wiki.python.org/moin/FrenchLanguage}{page sur le
  langage en français}.
\item
  La \href{https://docs.python.org/3/}{documentation originale} de
  Python 3 - donc, en anglais - est un très bon point d'entrée lorsqu'on
  cherche un sujet particulier, mais (beaucoup) trop abondante pour être
  lue d'un seul trait. Pour chercher de la documentation sur un module
  particulier, le plus simple est encore d'utiliser Google - ou votre
  moteur de recherche favori - qui vous redirigera, dans la grande
  majorité des cas, vers la page qui va bien dans, précisément, la
  documentation de Python.

  À titre d'exercice, cherchez la documentation du module
  \texttt{pathlib} en
  \href{https://www.google.fr/search?q=python+module+pathlib}{cherchant
  sur Google} les mots-clé \texttt{"python\ module\ pathlib"}.
\item
  J'aimerais vous signaler également une initiative pour
  \href{https://docs.python.org/fr/3/}{traduire la documentation
  officielle en français}.
\end{itemize}

    \hypertarget{historique-et-survol}{%
\subsubsection{Historique et survol}\label{historique-et-survol}}

    \begin{itemize}
\tightlist
\item
  La FAQ officielle de Python (en anglais) sur
  \href{https://docs.python.org/3/faq/design.html}{les choix de
  conception et l'historique du langage}.
\item
  L'article de Wikipédia (en anglais) sur
  l'\href{http://en.wikipedia.org/wiki/History_of_Python}{historique du
  langage}.
\item
  Sur Wikipédia, un article (en anglais) sur
  \href{http://en.wikipedia.org/wiki/Python_syntax_and_semantics}{la
  syntaxe et la sémantique de Python}.
\end{itemize}

    \hypertarget{un-peu-de-folklore}{%
\subsubsection{Un peu de folklore}\label{un-peu-de-folklore}}

    \begin{itemize}
\tightlist
\item
  Le \href{https://www.youtube.com/watch?v=YgtL4S7Hrwo}{discours de
  Guido van Rossum à PyCon 2016}.
\item
  Sur YouTube, le
  \href{https://www.youtube.com/watch?v=anwy2MPT5RE}{sketch des Monty
  Python} d'où proviennent les termes \texttt{spam}, \texttt{eggs} et
  autres \texttt{beans} que l'on utilise traditionnellement dans les
  exemples en Python plutôt que \texttt{foo} et \texttt{bar}.
\item
  L'\href{http://en.wikipedia.org/wiki/Spam_\%28Monty_Python\%29}{article
  Wikipédia correspondant}, qui cite le langage Python.
\end{itemize}

    \hypertarget{compluxe9ment---niveau-intermuxe9diaire}{%
\subsection{Complément - niveau
intermédiaire}\label{compluxe9ment---niveau-intermuxe9diaire}}

    \hypertarget{licence}{%
\subsubsection{Licence}\label{licence}}

    \begin{itemize}
\tightlist
\item
  La \href{https://docs.python.org/3/license.html}{licence d'utilisation
  est disponible ici}.
\item
  La page de la \href{https://www.python.org/psf/}{Python Software
  Foundation}, qui est une entité légale similaire à nos associations de
  1901, à but non lucratif~; elle possède les droits sur le langage.
\end{itemize}

    \hypertarget{le-processus-de-duxe9veloppement}{%
\subsubsection{Le processus de
développement}\label{le-processus-de-duxe9veloppement}}

    \begin{itemize}
\tightlist
\item
  Comment les choix d'évolution sont proposés et discutés,
  \href{http://en.wikipedia.org/wiki/Python_\%28programming_language\%29\#Development}{au
  travers des PEP (Python Enhancement Proposals) - sur wikipedia}
\item
  \href{http://legacy.python.org/dev/peps/pep-0001/}{Le premier PEP:
  PEP-001 donc} décrit en détail le cycle de vie des PEPs
\item
  Le \href{http://legacy.python.org/dev/peps/pep-0008}{PEP 008}, qui
  préconise un style de présentation (\emph{style guide})
\item
  L'\href{http://legacy.python.org/dev/peps/}{index de tous les PEPs}
\end{itemize}


    % Add a bibliography block to the postdoc
    
    
    

        \hypertarget{notebooks-jupyter-comme-support-de-cours}{%
\section{``Notebooks'' Jupyter comme support de
cours}\label{notebooks-jupyter-comme-support-de-cours}}

    Pour illustrer les vidéos du MOOC, nous avons choisi d'utiliser Jupyter
pour vous rédiger les documents ``mixtes'' contenant du texte et du code
Python, qu'on appelle des ``notebooks'', et dont le présent document est
un exemple.

Nous allons, dans la suite, utiliser du code Python, pourtant nous
n'avons pas encore abordé le langage. Pas d'inquiétude, ce code est
uniquement destiné à valider le fonctionnement des notebooks, et nous
n'utilisons que des choses très simples.

    \hypertarget{avertissement-ruxe9glages-du-navigateur}{%
\subsubsection{Avertissement: réglages du
navigateur}\label{avertissement-ruxe9glages-du-navigateur}}

    Avant toute chose, pour un bon fonctionnement des notebooks, on rappelle
qu'il est nécessaire d'avoir \textbf{autorisé} dans votre navigateur les
\textbf{cookies} en provenance du site Internet
\textbf{\texttt{nbhosting.inria.fr}}, qui héberge l'infrastructure qui
héberge tous les notebooks.

    \hypertarget{avantages-des-notebooks}{%
\subsubsection{Avantages des notebooks}\label{avantages-des-notebooks}}

    Comme vous le voyez, ce support permet un format plus lisible que des
commentaires dans un fichier de code.

    Nous attirons votre attention sur le fait que \textbf{les fragments de
code peuvent être évalués et modifiés}. Ainsi vous pouvez facilement
essayer des variantes autour du notebook original.

Notez bien également que le code Python est interprété \textbf{sur une
machine distante}, ce qui vous permet de faire vos premiers pas avant
même d'avoir procédé à l'installation de Python sur votre propre
ordinateur.

    \hypertarget{comment-utiliser-les-notebooks}{%
\subsubsection{Comment utiliser les
notebooks}\label{comment-utiliser-les-notebooks}}

    En haut du notebook, vous avez une barre de menu (sur fond bleu clair),
contenant~:

\begin{itemize}
	\item
un titre pour le notebook, avec un numéro de version~;
	\item
une barre de menus avec les entrées \texttt{File}, \texttt{Insert},
\texttt{Cell}, \texttt{Kernel};
	\item
et une barre de boutons qui sont des
raccourcis vers certains menus fréquemment utilisés. Si vous laissez
votre souris au dessus d'un bouton, un petit texte apparaît, indiquant à
quelle fonction correspond ce bouton.
\end{itemize}

Nous avons vu dans la vidéo qu'un notebook est constitué d'une suite de
cellules, soit textuelles, soit contenant du code. Les cellules de code
sont facilement reconnaissables, elles sont précédées de
\texttt{In\ {[}\ {]}:}. La cellule qui suit celle que vous êtes en train
de lire est une cellule de code.

Pour commencer, sélectionnez cette cellule de code avec votre souris, et
appuyez dans la barre de menu - en haut du notebook, donc - sur celui en
forme de flèche triangulaire vers la droite (Play)~: 

\includegraphics{medias/notebook-eval-button.png}\\
    \begin{Verbatim}[commandchars=\\\{\}]
{\color{incolor}In [{\color{incolor}1}]:} \PY{l+m+mi}{20} \PY{o}{*} \PY{l+m+mi}{30}
\end{Verbatim}


\begin{Verbatim}[commandchars=\\\{\}]
{\color{outcolor}Out[{\color{outcolor}1}]:} 600
\end{Verbatim}
            
    Comme vous le voyez, la cellule est ``exécutée'' (on dira plus
volontiers évaluée), et on passe à la cellule suivante.

Alternativement, vous pouvez simplement taper au clavier
\textbf{\emph{Shift+Enter}}, ou selon les claviers
\textbf{\emph{Maj-Entrée}}, pour obtenir le même effet. D'une manière
générale, il est important d'apprendre et d'utiliser les raccourcis
clavier, cela vous fera gagner beaucoup de temps par la suite.

    La façon habituelle d'\emph{exécuter} l'ensemble du notebook consiste~:
\begin{itemize}
	\item
	à sélectionner la première cellule,
	\item
	et à taper \textbf{\emph{Shift+Enter}} jusqu'à attendre la fin du notebook.
\end{itemize}

    Lorsqu'une cellule de code a été évaluée, Jupyter ajoute sous la cellule
\texttt{In} une cellule \texttt{Out} qui donne le résultat du fragment
Python, soit ci-dessus 600.\\

Jupyter ajoute également un nombre entre les crochets pour afficher, par
exemple ci-dessus, \texttt{In\ {[}1{]}:}. Ce nombre vous permet de
retrouver l'ordre dans lequel les cellules ont été évaluées.

    Vous pouvez naturellement modifier ces cellules de code pour faire des
essais~; ainsi vous pouvez vous servir du modèle ci-dessous pour
calculer la racine carrée de 3, ou essayer la fonction sur un nombre
négatif et voir comment est signalée l'erreur.

    \begin{Verbatim}[commandchars=\\\{\}]
{\color{incolor}In [{\color{incolor}2}]:} \PY{c+c1}{\PYZsh{} math.sqrt (pour square root) calcule la racine carrée}
        \PY{k+kn}{import} \PY{n+nn}{math}
        \PY{n}{math}\PY{o}{.}\PY{n}{sqrt}\PY{p}{(}\PY{l+m+mi}{2}\PY{p}{)}
\end{Verbatim}


\begin{Verbatim}[commandchars=\\\{\}]
{\color{outcolor}Out[{\color{outcolor}2}]:} 1.4142135623730951
\end{Verbatim}
            
    On peut également évaluer tout le notebook en une seule fois en
utilisant le menu \emph{Cell -\textgreater{} Run All}.

    \hypertarget{attention-uxe0-bien-uxe9valuer-les-cellules-dans-lordre}{%
\subsubsection{Attention à bien évaluer les cellules dans
l'ordre}\label{attention-uxe0-bien-uxe9valuer-les-cellules-dans-lordre}}

    Il est important que les cellules de code soient évaluées dans le bon
ordre. Si vous ne respectez pas l'ordre dans lequel les cellules de code
sont présentées, le résultat peut être inattendu.\\

En fait, évaluer un programme sous forme de notebook revient à le
découper en petits fragments, et si on exécute ces fragments dans le
désordre, on obtient naturellement un programme différent.

    On le voit sur cet exemple~:

    \begin{Verbatim}[commandchars=\\\{\}]
{\color{incolor}In [{\color{incolor}3}]:} \PY{n}{message} \PY{o}{=} \PY{l+s+s2}{\PYZdq{}}\PY{l+s+s2}{Il faut faire attention à l}\PY{l+s+s2}{\PYZsq{}}\PY{l+s+s2}{ordre dans lequel on évalue les notebooks}\PY{l+s+s2}{\PYZdq{}}
\end{Verbatim}


    \begin{Verbatim}[commandchars=\\\{\}]
{\color{incolor}In [{\color{incolor}4}]:} \PY{n+nb}{print}\PY{p}{(}\PY{n}{message}\PY{p}{)}
\end{Verbatim}


    \begin{Verbatim}[commandchars=\\\{\}]
Il faut faire attention à l'ordre dans lequel on évalue les notebooks

    \end{Verbatim}

    Si un peu plus loin dans le notebook on fait par exemple~:

    \begin{Verbatim}[commandchars=\\\{\}]
{\color{incolor}In [{\color{incolor}5}]:} \PY{c+c1}{\PYZsh{} ceci a pour effet d\PYZsq{}effacer la variable \PYZsq{}message\PYZsq{}}
        \PY{k}{del} \PY{n}{message}
\end{Verbatim}


    qui rend le symbole \texttt{message} indéfini, alors bien sûr on ne peut
plus évaluer la cellule qui fait \texttt{print} puisque la variable
\texttt{message} n'est plus connue de l'interpréteur.

    \hypertarget{ruxe9initialiser-linterpruxe9teur}{%
\subsubsection{Réinitialiser
l'interpréteur}\label{ruxe9initialiser-linterpruxe9teur}}

    Si vous faites trop de modifications, ou perdez le fil de ce que vous
avez évalué, il peut être utile de redémarrer votre interpréteur. Le
menu \emph{Kernel~→~Restart} vous permet de faire cela, un peu à la
manière de IDLE qui repart d'un interpréteur vierge lorsque vous
utilisez la fonction F5.

    Le menu \emph{Kernel~→~Interrupt} peut être quant à lui utilisé si votre
fragment prend trop longtemps à s'exécuter (par exemple vous avez écrit
une boucle dont la logique est cassée et qui ne termine pas).

    \hypertarget{vous-travaillez-sur-une-copie}{%
\subsubsection{Vous travaillez sur une
copie}\label{vous-travaillez-sur-une-copie}}

    Un des avantages principaux des notebooks est de vous permettre de
modifier le code que nous avons écrit, et de voir par vous-mêmes comment
se comporte le code modifié.\\

Pour cette raison, chaque élève dispose de sa \textbf{propre copie} de
chaque notebook, vous pouvez bien sûr apporter toutes les modifications
que vous souhaitez à vos notebooks sans affecter les autres étudiants.

    \hypertarget{revenir-uxe0-la-version-du-cours}{%
\subsubsection{Revenir à la version du
cours}\label{revenir-uxe0-la-version-du-cours}}

    Vous pouvez toujours revenir à la version ``du cours'' grâce au menu
\emph{File~→~Reset to original}.\\

    Attention, avec cette fonction vous restaurez \textbf{tout le notebook}
et donc \textbf{vous perdez vos modifications sur ce notebook}.

    \hypertarget{tuxe9luxe9charger-au-format-python}{%
\subsubsection{Télécharger au format
Python}\label{tuxe9luxe9charger-au-format-python}}

    Vous pouvez télécharger un notebook au format Python sur votre
ordinateur grâce au menu \emph{File~→~Download~as~→~Python}

    Les cellules de texte sont préservées dans le résultat sous forme de
commentaires Python.

    \hypertarget{partager-un-notebook-en-lecture-seule}{%
\subsubsection{Partager un notebook en lecture
seule}\label{partager-un-notebook-en-lecture-seule}}

    Enfin, avec le menu \emph{File~→~Share~static~version}, vous pouvez
publier une version en lecture seule de votre notebook~; vous obtenez
une URL que vous pouvez publier, par exemple pour demander de l'aide sur
le forum. Ainsi, les autres étudiants peuvent accéder en lecture seule à
votre code.\\

    Notez que lorsque vous utilisez cette fonction plusieurs fois, c'est
toujours la dernière version publiée que verront vos camarades, l'URL
utilisée reste toujours la même pour un étudiant et un notebook donné.

    \hypertarget{ajouter-des-cellules}{%
\subsubsection{Ajouter des cellules}\label{ajouter-des-cellules}}

    Vous pouvez ajouter une cellule n'importe où dans le document avec le
bouton \textbf{+} de la barre de boutons.\\

Aussi, lorsque vous arrivez à la fin du document, une nouvelle cellule
est créée chaque fois que vous évaluez la dernière cellule~; de cette
façon vous disposez d'un brouillon pour vos propres essais.\\

À vous de jouer.

        
    
    
    

    

    \hypertarget{modes-dexuxe9cution}{%
\section{Modes d'exécution}\label{modes-dexuxe9cution}}

    Nous avons donc à notre disposition plusieurs façons d'exécuter un
programme Python. Nous allons les étudier plus en détail~:

    \begin{longtable}[]{@{}ll@{}}
\toprule
Quoi & Avec quel outil\tabularnewline
\midrule
\endhead
fichier complet &
\texttt{python3\ \textless{}fichier\textgreater{}.py}\tabularnewline
ligne à ligne & \texttt{python3} en mode interactif\tabularnewline
~ & ou sous \texttt{ipython3}\tabularnewline
~ & ou avec IDLE\tabularnewline
par fragments & dans un notebook\tabularnewline
\bottomrule
\end{longtable}

    Pour cela nous allons voir le comportement d'un tout petit programme
Python lorsqu'on l'exécute sous ces différents environnements.

On veut surtout expliquer une petite différence quant au niveau de
détail de ce qui se trouve imprimé.

    Essentiellement, lorsqu'on utilise l'interpréteur en mode interactif -
ou sous IDLE - à chaque fois que l'on tape une ligne, le résultat est
\textbf{calculé} (on dit aussi \textbf{évalué}) puis \textbf{imprimé}.

Par contre, lorsqu'on écrit tout un programme, on ne peut plus imprimer
le résultat de toutes les lignes, cela produirait un flot d'impression
beaucoup trop important. Par conséquent, si vous ne déclenchez pas une
impression avec, par exemple, la fonction \texttt{print}, rien ne
s'affichera.

Enfin, en ce qui concerne le notebook, le comportement est un peu
hybride entre les deux, en ce sens que seul le \textbf{dernier résultat}
de la cellule est imprimé.

    \hypertarget{linterpruxe9teur-python-interactif}{%
\subsubsection{L'interpréteur Python
interactif}\label{linterpruxe9teur-python-interactif}}

    Le programme choisi est très simple, c'est le suivant~:

    \begin{Shaded}
\begin{Highlighting}[frame=lines,framerule=0.6mm,rulecolor=\color{asisframecolor}]
\DecValTok{10} \OperatorTok{*} \DecValTok{10}
\DecValTok{20} \OperatorTok{*} \DecValTok{20}
\DecValTok{30} \OperatorTok{*} \DecValTok{30}
\end{Highlighting}
\end{Shaded}

    Voici comment se comporte l'interpréteur interactif quand on lui soumet
ces instructions~:

    \begin{Shaded}
\begin{Highlighting}[frame=lines,framerule=0.6mm,rulecolor=\color{asisframecolor}]
\NormalTok{$ python3}
\NormalTok{Python }\DecValTok{3}\NormalTok{.}\FloatTok{5.1}\NormalTok{ (v3.}\FloatTok{5.1}\NormalTok{:37a07cee5969, Dec  }\DecValTok{5} \DecValTok{2015}\NormalTok{, }\DecValTok{21}\NormalTok{:}\DecValTok{12}\NormalTok{:}\DecValTok{44}\NormalTok{)}
\NormalTok{[GCC }\DecValTok{4}\NormalTok{.}\FloatTok{2.1}\NormalTok{ (Apple Inc. build }\DecValTok{5666}\NormalTok{) (dot }\DecValTok{3}\NormalTok{)] on darwin}
\NormalTok{Type }\StringTok{"help"}\NormalTok{, }\StringTok{"copyright"}\NormalTok{, }\StringTok{"credits"} \KeywordTok{or} \StringTok{"license"} \ControlFlowTok{for}\NormalTok{ more information.}
\OperatorTok{>>>} \DecValTok{10} \OperatorTok{*} \DecValTok{10}
\DecValTok{100}
\OperatorTok{>>>} \DecValTok{20} \OperatorTok{*} \DecValTok{20}
\DecValTok{400}
\OperatorTok{>>>} \DecValTok{30} \OperatorTok{*} \DecValTok{30}
\DecValTok{900}
\OperatorTok{>>>}\NormalTok{ exit()}
\NormalTok{$}
\end{Highlighting}
\end{Shaded}

    Notez que pour terminer la session, il nous faut ``sortir'' de
l'interpréteur en tapant \texttt{exit()}.

On peut aussi taper \texttt{Control-D} sous Linux ou macOS.

    Comme on le voit ici, l'interpréteur imprime \textbf{le résultat de
chaque ligne}. On voit bien apparaître toutes les valeurs calculées,
100, 400, puis enfin 900.

    \hypertarget{sous-forme-de-programme-constituuxe9}{%
\subsubsection{Sous forme de programme
constitué}\label{sous-forme-de-programme-constituuxe9}}

    Voyons à présent ce que donne cette même séquence de calculs dans un
programme complet. Pour cela, il nous faut tout d'abord fabriquer un
fichier avec un suffixe en \texttt{.py}, en utilisant par exemple un
éditeur de fichier. Le résultat doit ressembler à ceci~:

    \begin{Shaded}
\begin{Highlighting}[frame=lines,framerule=0.6mm,rulecolor=\color{asisframecolor}]
\NormalTok{$ }\FunctionTok{cat}\NormalTok{ foo.py}
\ExtensionTok{10}\NormalTok{ * 10}
\ExtensionTok{20}\NormalTok{ * 20}
\ExtensionTok{30}\NormalTok{ * 30}
\NormalTok{$}
\end{Highlighting}
\end{Shaded}

    Exécutons à présent ce programme~:

    \begin{Shaded}
\begin{Highlighting}[frame=lines,framerule=0.6mm,rulecolor=\color{asisframecolor}]
\NormalTok{$ }\ExtensionTok{python3}\NormalTok{ foo.py}
\NormalTok{$}
\end{Highlighting}
\end{Shaded}

    On constate donc que ce programme \textbf{ne fait rien~!} En tout cas,
selon toute apparence.

En réalité, les 3 valeurs 100, 400 et 900 sont bien calculées, mais
comme aucune instruction \texttt{print} n'est présente, rien n'est
imprimé et le programme se termine sans signe apparent d'avoir
réellement fonctionné.

Ce comportement peut paraître un peu déroutant au début, mais comme nous
l'avons mentionné c'est tout à fait délibéré. Un programme fonctionnel
faisant facilement plusieurs milliers de lignes, voire beaucoup plus, il
ne serait pas du tout réaliste que chaque ligne produise une impression,
comme c'est le cas en mode interactif.

    \hypertarget{dans-un-notebook}{%
\subsubsection{Dans un notebook}\label{dans-un-notebook}}

    Voici à présent le même programme dans un notebook~:

    \begin{Verbatim}[commandchars=\\\{\},frame=single,framerule=0.3mm,rulecolor=\color{cellframecolor}]
{\color{incolor}In [{\color{incolor}1}]:} \PY{l+m+mi}{10} \PY{o}{*} \PY{l+m+mi}{10}
        \PY{l+m+mi}{20} \PY{o}{*} \PY{l+m+mi}{20}
        \PY{l+m+mi}{30} \PY{o}{*} \PY{l+m+mi}{30}
\end{Verbatim}


\begin{Verbatim}[commandchars=\\\{\},frame=single,framerule=0.3mm,rulecolor=\color{cellframecolor}]
{\color{outcolor}Out[{\color{outcolor}1}]:} 900
\end{Verbatim}
            
    Lorsqu'on exécute cette cellule (rappel~: sélectionner la cellule, et
utiliser le bouton en forme de flèche vers la droite, ou entrer
\textbf{``Shift+Enter''} au clavier), on obtient une seule valeur dans
la rubrique \texttt{Out{[}{]}}, 900, qui correspond \textbf{au résultat
de la dernière ligne.}

    \hypertarget{utiliser-print}{%
\subsubsection{\texorpdfstring{Utiliser
\texttt{print}}{Utiliser print}}\label{utiliser-print}}

    Ainsi, pour afficher un résultat intermédiaire, on utilise l'instruction
\texttt{print}. Nous verrons cette instruction en détail dans les
semaines qui viennent, mais en guise d'introduction disons seulement que
c'est une fonction comme les autres en Python 3.

    \begin{Verbatim}[commandchars=\\\{\},frame=single,framerule=0.3mm,rulecolor=\color{cellframecolor}]
{\color{incolor}In [{\color{incolor}2}]:} \PY{n}{a} \PY{o}{=} \PY{l+m+mi}{10}
        \PY{n}{b} \PY{o}{=} \PY{l+m+mi}{20}
        
        \PY{n+nb}{print}\PY{p}{(}\PY{n}{a}\PY{p}{,} \PY{n}{b}\PY{p}{)}
\end{Verbatim}


    \begin{Verbatim}[commandchars=\\\{\},frame=single,framerule=0.3mm,rulecolor=\color{cellframecolor}]
10 20
\end{Verbatim}

    On peut naturellement mélanger des objets de plusieurs types, et donc
mélanger des chaînes de caractères et des nombres pour obtenir un
résultat un peu plus lisible. En effet, lorsque le programme devient
gros, il est important de savoir à quoi correspond une ligne dans le
flot de toutes les impressions. Aussi on préfèrera quelque chose comme~:

    \begin{Verbatim}[commandchars=\\\{\},frame=single,framerule=0.3mm,rulecolor=\color{cellframecolor}]
{\color{incolor}In [{\color{incolor}3}]:} \PY{n+nb}{print}\PY{p}{(}\PY{l+s+s2}{\PYZdq{}}\PY{l+s+s2}{a =}\PY{l+s+s2}{\PYZdq{}}\PY{p}{,} \PY{n}{a}\PY{p}{,} \PY{l+s+s2}{\PYZdq{}}\PY{l+s+s2}{et b =}\PY{l+s+s2}{\PYZdq{}}\PY{p}{,} \PY{n}{b}\PY{p}{)}
\end{Verbatim}


    \begin{Verbatim}[commandchars=\\\{\},frame=single,framerule=0.3mm,rulecolor=\color{cellframecolor}]
a = 10 et b = 20
\end{Verbatim}

    \begin{Verbatim}[commandchars=\\\{\},frame=single,framerule=0.3mm,rulecolor=\color{cellframecolor}]
{\color{incolor}In [{\color{incolor}4}]:} \PY{c+c1}{\PYZsh{} ou encore, équivalente mais avec un f\PYZhy{}string}
        \PY{n+nb}{print}\PY{p}{(}\PY{n}{f}\PY{l+s+s2}{\PYZdq{}}\PY{l+s+s2}{a = }\PY{l+s+si}{\PYZob{}a\PYZcb{}}\PY{l+s+s2}{ et b = }\PY{l+s+si}{\PYZob{}b\PYZcb{}}\PY{l+s+s2}{\PYZdq{}}\PY{p}{)}
\end{Verbatim}


    \begin{Verbatim}[commandchars=\\\{\},frame=single,framerule=0.3mm,rulecolor=\color{cellframecolor}]
a = 10 et b = 20
\end{Verbatim}

    Une pratique courante consiste d'ailleurs à utiliser les commentaires
pour laisser dans le code les instructions \texttt{print} qui
correspondent à du debug (c'est-à-dire qui ont pu être utiles lors de la
mise au point et qu'on veut pouvoir réactiver rapidement).

    \hypertarget{utiliser-print-pour-sous-titrer-une-affectation}{%
\subsubsection{\texorpdfstring{Utiliser \texttt{print} pour
``sous-titrer'' une
affectation}{Utiliser print pour ``sous-titrer'' une affectation}}\label{utiliser-print-pour-sous-titrer-une-affectation}}

    Remarquons enfin que l'affectation à une variable ne retourne aucun
résultat.

C'est-à-dire, en pratique, que si on écrit~:

    \begin{Verbatim}[commandchars=\\\{\},frame=single,framerule=0.3mm,rulecolor=\color{cellframecolor}]
{\color{incolor}In [{\color{incolor}5}]:} \PY{n}{a} \PY{o}{=} \PY{l+m+mi}{100}
\end{Verbatim}


    même une fois l'expression évaluée par l'interpréteur, aucune ligne
\texttt{Out{[}{]}} n'est ajoutée.

    C'est pourquoi, il nous arrivera parfois d'écrire, notamment lorsque
l'expression est complexe et pour rendre explicite la valeur qui vient
d'être affectée~:

    \begin{Verbatim}[commandchars=\\\{\},frame=single,framerule=0.3mm,rulecolor=\color{cellframecolor}]
{\color{incolor}In [{\color{incolor}6}]:} \PY{n}{a} \PY{o}{=} \PY{l+m+mi}{100}\PY{p}{;} \PY{n+nb}{print}\PY{p}{(}\PY{n}{a}\PY{p}{)}
\end{Verbatim}


    \begin{Verbatim}[commandchars=\\\{\},frame=single,framerule=0.3mm,rulecolor=\color{cellframecolor}]
100
\end{Verbatim}

    Notez bien que cette technique est uniquement pédagogique, et n'a
absolument aucun autre intérêt dans la pratique~; il n'est \textbf{pas
recommandé} de l'utiliser en dehors de ce contexte.


    % Add a bibliography block to the postdoc
    
    
    

        
    
    
    

    

    \hypertarget{la-suite-de-fibonacci}{%
\section{La suite de Fibonacci}\label{la-suite-de-fibonacci}}

    \hypertarget{compluxe9ment---niveau-basique}{%
\subsection{Complément - niveau
basique}\label{compluxe9ment---niveau-basique}}

    Voici un premier exemple de code qui tourne.

Nous allons commencer par le faire tourner dans ce notebook. Nous
verrons en fin de séance comment le faire fonctionner localement sur
votre ordinateur.

    Le but de ce programme est de calculer la
\href{https://fr.wikipedia.org/wiki/Suite_de_Fibonacci}{suite de
Fibonacci}, qui est définie comme ceci~:

\begin{itemize}
\tightlist
\item
  \(u_0 = 1\)
\item
  \(u_1 = 1\)
\item
  \(\forall n >= 2, u_n = u_{n-1} + u_{n-2}\)
\end{itemize}

Ce qui donne pour les premières valeurs~:

    \begin{longtable}[]{@{}ll@{}}
\toprule
n & fibonacci(n)\tabularnewline
\midrule
\endhead
0 & 1\tabularnewline
1 & 1\tabularnewline
2 & 2\tabularnewline
3 & 3\tabularnewline
4 & 5\tabularnewline
5 & 8\tabularnewline
6 & 13\tabularnewline
\bottomrule
\end{longtable}

    On commence par définir la fonction \texttt{fibonacci} comme il suit.
Naturellement vous n'avez pas encore tout le bagage pour lire ce code,
ne vous inquiétez pas, nous allons vous expliquer tout ça dans les
prochaines semaines. Le but est uniquement de vous montrer un
fonctionnement de l'interpréteur Python et de IDLE.

    \begin{Verbatim}[commandchars=\\\{\},frame=single,framerule=0.3mm,rulecolor=\color{cellframecolor}]
{\color{incolor}In [{\color{incolor}1}]:} \PY{k}{def} \PY{n+nf}{fibonacci}\PY{p}{(}\PY{n}{n}\PY{p}{)}\PY{p}{:}
            \PY{l+s+s2}{\PYZdq{}}\PY{l+s+s2}{retourne le nombre de fibonacci pour l}\PY{l+s+s2}{\PYZsq{}}\PY{l+s+s2}{entier n}\PY{l+s+s2}{\PYZdq{}}
            \PY{c+c1}{\PYZsh{} pour les petites valeurs de n il n\PYZsq{}y a rien à calculer}
            \PY{k}{if} \PY{n}{n} \PY{o}{\PYZlt{}}\PY{o}{=} \PY{l+m+mi}{1}\PY{p}{:}
                \PY{k}{return} \PY{l+m+mi}{1}
            \PY{c+c1}{\PYZsh{} sinon on initialise f1 pour n\PYZhy{}1 et f2 pour n\PYZhy{}2}
            \PY{n}{f2}\PY{p}{,} \PY{n}{f1} \PY{o}{=} \PY{l+m+mi}{1}\PY{p}{,} \PY{l+m+mi}{1}
            \PY{c+c1}{\PYZsh{} et on itère n\PYZhy{}1 fois pour additionner}
            \PY{k}{for} \PY{n}{i} \PY{o+ow}{in} \PY{n+nb}{range}\PY{p}{(}\PY{l+m+mi}{2}\PY{p}{,} \PY{n}{n} \PY{o}{+} \PY{l+m+mi}{1}\PY{p}{)}\PY{p}{:}
                \PY{n}{f2}\PY{p}{,} \PY{n}{f1} \PY{o}{=} \PY{n}{f1}\PY{p}{,} \PY{n}{f1} \PY{o}{+} \PY{n}{f2}
        \PY{c+c1}{\PYZsh{}        print(i, f2, f1)}
            \PY{c+c1}{\PYZsh{} le résultat est dans f1}
            \PY{k}{return} \PY{n}{f1}
\end{Verbatim}


    Pour en faire un programme utilisable on va demander à l'utilisateur de
rentrer un nombre~; il faut le convertir en entier car \texttt{input}
renvoie une chaîne de caractères~:

    \begin{Verbatim}[commandchars=\\\{\},frame=single,framerule=0.3mm,rulecolor=\color{cellframecolor}]
{\color{incolor}In [{\color{incolor}2}]:} \PY{n}{entier} \PY{o}{=} \PY{n+nb}{int}\PY{p}{(}\PY{n+nb}{input}\PY{p}{(}\PY{l+s+s2}{\PYZdq{}}\PY{l+s+s2}{Entrer un entier }\PY{l+s+s2}{\PYZdq{}}\PY{p}{)}\PY{p}{)}
        
        \PY{c+c1}{\PYZsh{} NOTE:}
        \PY{c+c1}{\PYZsh{} auto\PYZhy{}exec\PYZhy{}for\PYZhy{}latex has used instead:}
        \PY{c+c1}{\PYZsh{}\PYZsh{}\PYZsh{}\PYZsh{}\PYZsh{}\PYZsh{}\PYZsh{}\PYZsh{}\PYZsh{}\PYZsh{}}
        \PY{n}{entier} \PY{o}{=} \PY{l+m+mi}{12}
        \PY{c+c1}{\PYZsh{}\PYZsh{}\PYZsh{}\PYZsh{}\PYZsh{}\PYZsh{}\PYZsh{}\PYZsh{}\PYZsh{}\PYZsh{}}
\end{Verbatim}


    On imprime le résultat~:

    \begin{Verbatim}[commandchars=\\\{\},frame=single,framerule=0.3mm,rulecolor=\color{cellframecolor}]
{\color{incolor}In [{\color{incolor}3}]:} \PY{n+nb}{print}\PY{p}{(}\PY{n}{f}\PY{l+s+s2}{\PYZdq{}}\PY{l+s+s2}{fibonacci(}\PY{l+s+si}{\PYZob{}entier\PYZcb{}}\PY{l+s+s2}{) = }\PY{l+s+s2}{\PYZob{}}\PY{l+s+s2}{fibonacci(entier)\PYZcb{}}\PY{l+s+s2}{\PYZdq{}}\PY{p}{)}
\end{Verbatim}


    \begin{Verbatim}[commandchars=\\\{\},frame=single,framerule=0.3mm,rulecolor=\color{cellframecolor}]
fibonacci(12) = 233
\end{Verbatim}

    \hypertarget{exercice}{%
\subsubsection{Exercice}\label{exercice}}

    Vous pouvez donc à présent~:

\begin{itemize}
\tightlist
\item
  exécuter le code dans ce notebook
\item
  télécharger ce code sur votre disque comme un fichier
  \texttt{fibonacci\_prompt.py}

  \begin{itemize}
  \tightlist
  \item
    utiliser pour cela le menu
    \emph{``File~-\textgreater{}~Download~as~-\textgreater{}~Python''}
  \item
    et \textbf{renommer le fichier obtenu} au besoin
  \end{itemize}
\item
  l'exécuter sous IDLE
\item
  le modifier, par exemple pour afficher les résultats intermédiaires

  \begin{itemize}
  \tightlist
  \item
    on a laissé exprès une fonction \texttt{print} en commentaire que
    vous pouvez réactiver simplement
  \end{itemize}
\item
  l'exécuter avec l'interpréteur Python comme ceci~:
\end{itemize}

\begin{Shaded}
\begin{Highlighting}[frame=lines,framerule=0.6mm,rulecolor=\color{asisframecolor}]
\NormalTok{$ }\ExtensionTok{python3}\NormalTok{ fibonacci_prompt.py}
\end{Highlighting}
\end{Shaded}

    Ce code est volontairement simple et peu robuste pour ne pas l'alourdir.
Par exemple, ce programme se comporte mal si vous entrez un entier
négatif.

    Nous allons voir tout de suite une version légèrement différente qui va
vous permettre de donner la valeur d'entrée sur la ligne de commande.


    % Add a bibliography block to the postdoc
    
    
    

        \hypertarget{la-suite-de-fibonacci-suite}{%
\section{La suite de Fibonacci
(suite)}\label{la-suite-de-fibonacci-suite}}

    \hypertarget{compluxe9ment---niveau-intermuxe9diaire}{%
\subsection{Complément - niveau
intermédiaire}\label{compluxe9ment---niveau-intermuxe9diaire}}

    Nous reprenons le cas de la fonction \texttt{fibonacci} que nous avons
déjà vue, mais cette fois nous voulons que l'utilisateur puisse indiquer
l'entier en entrée de l'algorithme, non plus en répondant à une
question, mais sur la ligne de commande, c'est-à-dire en tapant~:

\begin{verbatim}
$ python3 fibonacci.py 12
\end{verbatim}

    \textbf{Avertissement~:}\\

Attention, cette version-ci \textbf{ne fonctionne pas dans ce notebook},
justement car on n'a pas de moyen dans un notebook d'invoquer un
programme en lui passant des arguments de cette façon. Ce notebook est
rédigé pour vous permettre de vous entraîner avec la fonction de
téléchargement au format Python, qu'on a vue dans la vidéo, et de faire
tourner ce programme sur votre propre ordinateur.

    \hypertarget{le-module-argparse}{%
\subsubsection{\texorpdfstring{Le module
\texttt{argparse}}{Le module argparse}}\label{le-module-argparse}}

    Cette fois nous importons le module \texttt{argparse}, c'est lui qui va
nous permettre d'interpréter les arguments passés sur la ligne de
commande.

    \begin{Verbatim}[commandchars=\\\{\}]
{\color{incolor}In [{\color{incolor} }]:} \PY{k+kn}{from} \PY{n+nn}{argparse} \PY{k}{import} \PY{n}{ArgumentParser}
\end{Verbatim}


    Puis nous répétons la fonction \texttt{fibonacci}~:

    \begin{Verbatim}[commandchars=\\\{\}]
{\color{incolor}In [{\color{incolor} }]:} \PY{k}{def} \PY{n+nf}{fibonacci}\PY{p}{(}\PY{n}{n}\PY{p}{)}\PY{p}{:}
            \PY{l+s+s2}{\PYZdq{}}\PY{l+s+s2}{retourne le nombre de fibonacci pour l}\PY{l+s+s2}{\PYZsq{}}\PY{l+s+s2}{entier n}\PY{l+s+s2}{\PYZdq{}}
            \PY{c+c1}{\PYZsh{} pour les petites valeurs de n il n\PYZsq{}y a rien à calculer}
            \PY{k}{if} \PY{n}{n} \PY{o}{\PYZlt{}}\PY{o}{=} \PY{l+m+mi}{1}\PY{p}{:}
                \PY{k}{return} \PY{l+m+mi}{1}
            \PY{c+c1}{\PYZsh{} sinon on initialise f1 pour n\PYZhy{}1 et f2 pour n\PYZhy{}2}
            \PY{n}{f2}\PY{p}{,} \PY{n}{f1} \PY{o}{=} \PY{l+m+mi}{1}\PY{p}{,} \PY{l+m+mi}{1}
            \PY{c+c1}{\PYZsh{} et on itère n\PYZhy{}1 fois pour additionner}
            \PY{k}{for} \PY{n}{i} \PY{o+ow}{in} \PY{n+nb}{range}\PY{p}{(}\PY{l+m+mi}{2}\PY{p}{,} \PY{n}{n} \PY{o}{+} \PY{l+m+mi}{1}\PY{p}{)}\PY{p}{:}
                \PY{n}{f2}\PY{p}{,} \PY{n}{f1} \PY{o}{=} \PY{n}{f1}\PY{p}{,} \PY{n}{f1} \PY{o}{+} \PY{n}{f2}
        \PY{c+c1}{\PYZsh{}        print(i, f2, f1)}
            \PY{c+c1}{\PYZsh{} le résultat est dans f1}
            \PY{k}{return} \PY{n}{f1}
\end{Verbatim}


    \emph{Remarque~:}

Certains d'entre vous auront évidemment remarqué que l'on aurait pu
éviter de copier-coller la fonction \texttt{fibonacci} comme cela~;
c'est à ça que servent les modules, mais nous n'en sommes pas là.

    \hypertarget{un-objet-parser}{%
\subsubsection{\texorpdfstring{Un objet
\texttt{parser}}{Un objet parser}}\label{un-objet-parser}}

    À présent, nous utilisons le module \texttt{argparse}, pour lui dire
qu'on attend exactement un argument sur la ligne de commande, et qu'il
doit être un entier. Ici encore, ne vous inquiétez pas si vous ne
comprenez pas tout le code. L'objectif est de vous donner un morceau de
code utilisable tout de suite, pour jouer avec votre interpréteur
Python.

    \begin{Verbatim}[commandchars=\\\{\}]
{\color{incolor}In [{\color{incolor} }]:} \PY{c+c1}{\PYZsh{} à nouveau : ceci n\PYZsq{}est pas conçu pour être exécuté dans le notebook !}
        \PY{n}{parser} \PY{o}{=} \PY{n}{ArgumentParser}\PY{p}{(}\PY{p}{)}
        \PY{n}{parser}\PY{o}{.}\PY{n}{add\PYZus{}argument}\PY{p}{(}\PY{n}{dest}\PY{o}{=}\PY{l+s+s2}{\PYZdq{}}\PY{l+s+s2}{entier}\PY{l+s+s2}{\PYZdq{}}\PY{p}{,} \PY{n+nb}{type}\PY{o}{=}\PY{n+nb}{int}\PY{p}{,}
                            \PY{n}{help}\PY{o}{=}\PY{l+s+s2}{\PYZdq{}}\PY{l+s+s2}{entier d}\PY{l+s+s2}{\PYZsq{}}\PY{l+s+s2}{entrée}\PY{l+s+s2}{\PYZdq{}}\PY{p}{)}
        \PY{n}{input\PYZus{}args} \PY{o}{=} \PY{n}{parser}\PY{o}{.}\PY{n}{parse\PYZus{}args}\PY{p}{(}\PY{p}{)}
        \PY{n}{entier} \PY{o}{=} \PY{n}{input\PYZus{}args}\PY{o}{.}\PY{n}{entier}
\end{Verbatim}


    Nous pouvons à présent afficher le résultat~:

    \begin{Verbatim}[commandchars=\\\{\}]
{\color{incolor}In [{\color{incolor} }]:} \PY{n+nb}{print}\PY{p}{(}\PY{n}{f}\PY{l+s+s2}{\PYZdq{}}\PY{l+s+s2}{fibonacci(}\PY{l+s+si}{\PYZob{}entier\PYZcb{}}\PY{l+s+s2}{) = }\PY{l+s+s2}{\PYZob{}}\PY{l+s+s2}{fibonacci(entier)\PYZcb{}}\PY{l+s+s2}{\PYZdq{}}\PY{p}{)}
\end{Verbatim}


    Vous pouvez donc à présent~:

\begin{itemize}
\item
  télécharger ce code sur votre disque comme un fichier
  \texttt{fibonacci.py} en utilisant le menu
  \emph{``File~-\textgreater{}~Download as~-\textgreater{}~Python''}
\item
  l'exécuter avec simplement l'interpréteur Python comme ceci~:

  \$ python3 fibonacci.py 56
\end{itemize}
        \hypertarget{la-ligne-shebang}{%
\section{\texorpdfstring{La ligne
\emph{shebang}}{La ligne shebang}}\label{la-ligne-shebang}}

    \begin{verbatim}
#!/usr/bin/env python3
\end{verbatim}

    \hypertarget{compluxe9ment---niveau-avancuxe9}{%
\subsection{Complément - niveau
avancé}\label{compluxe9ment---niveau-avancuxe9}}

    Ce complément est uniquement valable pour macOS et Linux.

    \hypertarget{le-besoin}{%
\subsubsection{Le besoin}\label{le-besoin}}

    Nous avons vu dans la vidéo que, pour lancer un programme Python, on
fait depuis le terminal~:

    \begin{verbatim}
$ python3 mon_module.py
\end{verbatim}

    Lorsqu'il s'agit d'un programme que l'on utilise fréquemment, on n'est
pas forcément dans le répertoire où se trouve le programme Python.
Aussi, dans ce cas, on peut utiliser un chemin ``absolu'', c'est-à-dire
à partir de la racine des noms de fichiers, comme par exemple~:

    \begin{verbatim}
$ python3 /le/chemin/jusqu/a/mon_module.py
\end{verbatim}

    Sauf que c'est assez malcommode, et cela devient vite pénible à la
longue.

    \hypertarget{la-solution}{%
\subsubsection{La solution}\label{la-solution}}

    Sur Linux et macOS, il existe une astuce utile pour simplifier cela.
Voyons comment s'y prendre, avec par exemple le programme
\texttt{fibonacci.py} que vous pouvez
\href{data/fibonacci.py}{télécharger ici} (nous avons vu ce code en
détail dans les deux compléments précédents). Commencez par sauver ce
code sur votre ordinateur dans un fichier qui s'appelle, bien entendu,
\texttt{fibonacci.py}.\\

    On commence par éditer le tout début du fichier pour lui ajouter une
\textbf{première ligne}~:

    \begin{verbatim}
#!/usr/bin/env python3

## La suite de Fibonacci (Suite)
...etc...
\end{verbatim}

    Cette première ligne s'appelle un
\href{http://en.wikipedia.org/wiki/Shebang_\%28Unix\%29}{Shebang} dans
le jargon Unix. Unix stipule que le Shebang doit être en
\textbf{première position} dans le fichier.\\

    Ensuite on rajoute au fichier, depuis le terminal, le caractère
exécutable comme ceci~:

    \begin{verbatim}
$ pwd
/le/chemin/jusqu/a/
\end{verbatim}

    \begin{verbatim}
$ chmod +x fibonacci.py
\end{verbatim}

    À partir de là, vous pouvez utiliser le fichier \texttt{fibonacci.py}
comme une commande, sans avoir à mentionner \texttt{python3}, qui sera
invoqué au travers du shebang~:

    \begin{verbatim}
$ /le/chemin/jusqu/a/fibonacci.py 20
fibonacci(20) = 10946
\end{verbatim}

    Et donc vous pouvez aussi le déplacer dans un répertoire qui est dans
votre variable \texttt{PATH}; de cette façon vous les rendez ainsi
accessible à partir n'importe quel répertoire en faisant simplement~:

    \begin{verbatim}
$ export PATH=/le/chemin/jusqu/a:$PATH
\end{verbatim}

    \begin{verbatim}
$ cd /tmp
$ fibonacci.py 20
fibonacci(20) = 10946
\end{verbatim}

    \textbf{Remarque~:} tout ceci fonctionne très bien tant que votre point
d'entrée - ici \texttt{fibonacci.py} - n'utilise que des modules
standards. Dans le cas où le point d'entrée vient avec au moins un
module, il est également nécessaire d'installer ces modules quelque
part, et d'indiquer au point d'entrée comment les trouver, nous y
reviendrons dans la semaine où nous parlerons des modules.
        
    
    
    

    

    \hypertarget{dessiner-un-carruxe9}{%
\section{Dessiner un carré}\label{dessiner-un-carruxe9}}

    \hypertarget{exercice---niveau-intermuxe9diaire}{%
\subsection{Exercice - niveau
intermédiaire}\label{exercice---niveau-intermuxe9diaire}}

    Voici un tout petit programme qui dessine un carré.

    Il utilise le module \texttt{turtle}, conçu précisément à des fins
pédagogiques. Pour des raisons techniques, le module \texttt{turtle}
n'est \textbf{pas disponible} au travers de la plateforme FUN.

    \textbf{Il est donc inutile d'essayer d'exécuter ce programme depuis le
notebook}. L'objectif de cet exercice est plutôt de vous entraîner à
télécharger ce programme en utilisant le menu
\emph{``File~-\textgreater{}~Download~as~-\textgreater{}~Python''}, puis
à le charger dans votre IDLE pour l'exécuter sur votre machine.

    \textbf{Attention} également à sauver le programme téléchargé
\textbf{sous un autre nom} que \texttt{turtle.py}, car sinon vous allez
empêcher Python de trouver le module standard \texttt{turtle}~;
appelez-le par exemple \texttt{turtle\_basic.py}.

    \begin{Verbatim}[commandchars=\\\{\},frame=single,framerule=0.3mm,rulecolor=\color{cellframecolor}]
{\color{incolor}In [{\color{incolor}1}]:} \PY{c+c1}{\PYZsh{} on a besoin du module turtle}
        \PY{k+kn}{import} \PY{n+nn}{turtle}
\end{Verbatim}


    On commence par définir une fonction qui dessine un carré de côté
\texttt{length}~:

    \begin{Verbatim}[commandchars=\\\{\},frame=single,framerule=0.3mm,rulecolor=\color{cellframecolor}]
{\color{incolor}In [{\color{incolor}2}]:} \PY{k}{def} \PY{n+nf}{square}\PY{p}{(}\PY{n}{length}\PY{p}{)}\PY{p}{:}
            \PY{l+s+s2}{\PYZdq{}}\PY{l+s+s2}{have the turtle draw a square of side \PYZlt{}length\PYZgt{}}\PY{l+s+s2}{\PYZdq{}}
            \PY{k}{for} \PY{n}{side} \PY{o+ow}{in} \PY{n+nb}{range}\PY{p}{(}\PY{l+m+mi}{4}\PY{p}{)}\PY{p}{:}
                \PY{n}{turtle}\PY{o}{.}\PY{n}{forward}\PY{p}{(}\PY{n}{length}\PY{p}{)}
                \PY{n}{turtle}\PY{o}{.}\PY{n}{left}\PY{p}{(}\PY{l+m+mi}{90}\PY{p}{)}
\end{Verbatim}


    Maintenant on commence par initialiser la tortue~:

    \begin{Verbatim}[commandchars=\\\{\},frame=single,framerule=0.3mm,rulecolor=\color{cellframecolor}]
{\color{incolor}In [{\color{incolor}3}]:} \PY{n}{turtle}\PY{o}{.}\PY{n}{reset}\PY{p}{(}\PY{p}{)}
\end{Verbatim}


    On peut alors dessiner notre carré~:

    \begin{Verbatim}[commandchars=\\\{\},frame=single,framerule=0.3mm,rulecolor=\color{cellframecolor}]
{\color{incolor}In [{\color{incolor} }]:} \PY{c+c1}{\PYZsh{} NOTE}
        \PY{c+c1}{\PYZsh{} auto\PYZhy{}exec\PYZhy{}for\PYZhy{}latex has skipped execution of this cell}
        
        \PY{n}{square}\PY{p}{(}\PY{l+m+mi}{200}\PY{p}{)}
\end{Verbatim}


    Et pour finir on attend que l'utilisateur clique dans la fenêtre de la
tortue, et alors on termine~:

    \begin{Verbatim}[commandchars=\\\{\},frame=single,framerule=0.3mm,rulecolor=\color{cellframecolor}]
{\color{incolor}In [{\color{incolor} }]:} \PY{c+c1}{\PYZsh{} NOTE}
        \PY{c+c1}{\PYZsh{} auto\PYZhy{}exec\PYZhy{}for\PYZhy{}latex has skipped execution of this cell}
        
        \PY{n}{turtle}\PY{o}{.}\PY{n}{exitonclick}\PY{p}{(}\PY{p}{)}
\end{Verbatim}


    \hypertarget{exercice---niveau-avancuxe9}{%
\subsection{Exercice - niveau
avancé}\label{exercice---niveau-avancuxe9}}

    Naturellement vous pouvez vous amuser à modifier ce code pour dessiner
des choses un peu plus amusantes.

Dans ce cas, commencez par chercher ``\emph{module python turtle}'' dans
votre moteur de recherche favori, pour localiser la documentation du
module
\href{https://docs.python.org/3/library/turtle.html}{\texttt{turtle}}.

Vous trouverez quelques exemples pour commencer ici~:

\begin{itemize}
\tightlist
\item
  \href{media/turtle_multi_squares.py}{turtle\_multi\_squares.py} pour
  dessiner des carrés à l'emplacement de la souris en utilisant
  plusieurs tortues~;
\item
  \href{media/turtle_fractal.py}{turtle\_fractal.py} pour dessiner une
  fractale simple~;
\item
  \href{media/turtle_fractal_reglable.py}{turtle\_fractal\_reglable.py}
  une variation sur la fractale, plus paramétrable.
\end{itemize}


    % Add a bibliography block to the postdoc
    
    
    

        \hypertarget{noms-de-variables}{%
\section{Noms de variables}\label{noms-de-variables}}

    \hypertarget{compluxe9ment---niveau-basique}{%
\subsection{Complément - niveau
basique}\label{compluxe9ment---niveau-basique}}

    Revenons sur les noms de variables autorisés ou non.\\

    Les noms les plus simples sont constitués de lettres. Par exemple~:

    \begin{Verbatim}[commandchars=\\\{\}]
{\color{incolor}In [{\color{incolor}1}]:} \PY{n}{factoriel} \PY{o}{=} \PY{l+m+mi}{1}
\end{Verbatim}


    On peut utiliser aussi les majuscules, mais attention cela définit une
variable différente. Ainsi~:

    \begin{Verbatim}[commandchars=\\\{\}]
{\color{incolor}In [{\color{incolor}2}]:} \PY{n}{Factoriel} \PY{o}{=} \PY{l+m+mi}{100}
        \PY{n}{factoriel} \PY{o}{==} \PY{n}{Factoriel}
\end{Verbatim}


\begin{Verbatim}[commandchars=\\\{\}]
{\color{outcolor}Out[{\color{outcolor}2}]:} False
\end{Verbatim}
            
    Le signe \texttt{==} permet de tester si deux variables ont la même
valeur. Si les variables ont la même valeur, le test retournera
\texttt{True}, et \texttt{False} sinon. On y reviendra bien entendu.

    \hypertarget{conventions-habituelles}{%
\subsubsection{Conventions habituelles}\label{conventions-habituelles}}

    En règle générale, on utilise \textbf{uniquement des minuscules} pour
désigner les variables simples (ainsi d'ailleurs que pour les noms de
fonctions), les majuscules sont réservées en principe pour d'autres
sortes de variables, comme les noms de classe, que nous verrons
ultérieurement.\\

Notons qu'il s'agit uniquement d'une convention, ceci n'est pas imposé
par le langage lui-même.\\

    Pour des raisons de lisibilité, il est également possible d'utiliser le
tiret bas \texttt{\_} dans les noms de variables. On préfèrera ainsi~:

    \begin{Verbatim}[commandchars=\\\{\}]
{\color{incolor}In [{\color{incolor}3}]:} \PY{n}{age\PYZus{}moyen} \PY{o}{=} \PY{l+m+mi}{75} \PY{c+c1}{\PYZsh{} oui}
\end{Verbatim}


    plutôt que ceci (bien qu'autorisé par le langage)~:

    \begin{Verbatim}[commandchars=\\\{\}]
{\color{incolor}In [{\color{incolor}4}]:} \PY{n}{AgeMoyen} \PY{o}{=} \PY{l+m+mi}{75} \PY{c+c1}{\PYZsh{} autorisé, mais non}
\end{Verbatim}


    On peut également utiliser des chiffres dans les noms de variables comme
par exemple~:

    \begin{Verbatim}[commandchars=\\\{\}]
{\color{incolor}In [{\color{incolor}5}]:} \PY{n}{age\PYZus{}moyen\PYZus{}dept75} \PY{o}{=} \PY{l+m+mi}{80}
\end{Verbatim}


    avec la restriction toutefois que le premier caractère ne peut pas être
un chiffre, cette affectation est donc refusée~:

    \begin{Verbatim}[commandchars=\\\{\}]
{\color{incolor}In [{\color{incolor}6}]:} \PY{l+m+mi}{75}\PY{n}{\PYZus{}age\PYZus{}moyen} \PY{o}{=} \PY{l+m+mi}{80} \PY{c+c1}{\PYZsh{} erreur de syntaxe}
\end{Verbatim}


    \begin{Verbatim}[commandchars=\\\{\}]

          File "<ipython-input-6-823fed77034a>", line 1
        75\_age\_moyen = 80 \# erreur de syntaxe
          \^{}
    SyntaxError: invalid token


    \end{Verbatim}

    \hypertarget{le-tiret-bas-comme-premier-caractuxe8re}{%
\subsubsection{Le tiret bas comme premier
caractère}\label{le-tiret-bas-comme-premier-caractuxe8re}}

    Il est par contre, possible de faire commencer un nom de variable par un
tiret bas comme premier caractère~; toutefois, à ce stade, nous vous
déconseillons d'utiliser cette pratique qui est réservée à des
conventions de nommage bien spécifiques.

    \begin{Verbatim}[commandchars=\\\{\}]
{\color{incolor}In [{\color{incolor}7}]:} \PY{n}{\PYZus{}autorise\PYZus{}mais\PYZus{}deconseille} \PY{o}{=} \PY{l+s+s1}{\PYZsq{}}\PY{l+s+s1}{Voir le PEP 008}\PY{l+s+s1}{\PYZsq{}}
\end{Verbatim}


    Et en tout cas, il est \textbf{fortement déconseillé} d'utiliser des
noms de la forme \texttt{\_\_variable\_\_} qui sont réservés au langage.
Nous reviendrons sur ce point dans le futur, mais regardez par exemple
cette variable que nous n'avons définie nulle part mais qui pourtant
existe bel et bien~:

    \begin{Verbatim}[commandchars=\\\{\}]
{\color{incolor}In [{\color{incolor}8}]:} \PY{n+nv+vm}{\PYZus{}\PYZus{}name\PYZus{}\PYZus{}}  \PY{c+c1}{\PYZsh{} ne définissez pas vous\PYZhy{}même de variables de ce genre}
\end{Verbatim}


\begin{Verbatim}[commandchars=\\\{\}]
{\color{outcolor}Out[{\color{outcolor}8}]:} '\_\_main\_\_'
\end{Verbatim}
            
    \hypertarget{ponctuation}{%
\subsubsection{Ponctuation}\label{ponctuation}}

    Dans la plage des caractères ASCII, il n'est \textbf{pas possible}
d'utiliser d'autres caractères que les caractères alphanumériques et le
tiret bas. Notamment le tiret haut \texttt{-} est interprété comme
l'opération de soustraction. Attention donc à cette erreur fréquente~:

    \begin{Verbatim}[commandchars=\\\{\}]
{\color{incolor}In [{\color{incolor}9}]:} \PY{n}{age}\PY{o}{\PYZhy{}}\PY{n}{moyen} \PY{o}{=} \PY{l+m+mi}{75}  \PY{c+c1}{\PYZsh{} erreur : en fait python l\PYZsq{}interprète comme \PYZsq{}age \PYZhy{} moyen = 75\PYZsq{}}
\end{Verbatim}


    \begin{Verbatim}[commandchars=\\\{\}]

          File "<ipython-input-9-78de3a1bfc60>", line 1
        age-moyen = 75  \# erreur : en fait python l'interprète comme 'age - moyen = 75'
                                                                                       \^{}
    SyntaxError: can't assign to operator


    \end{Verbatim}

    \hypertarget{caractuxe8res-exotiques}{%
\subsubsection{Caractères exotiques}\label{caractuxe8res-exotiques}}

    En Python 3, il est maintenant aussi possible d'utiliser des caractères
Unicode dans les identificateurs~:

    \begin{Verbatim}[commandchars=\\\{\}]
{\color{incolor}In [{\color{incolor}10}]:} \PY{c+c1}{\PYZsh{} les caractères accentués sont permis}
         \PY{n}{nom\PYZus{}élève} \PY{o}{=} \PY{l+s+s2}{\PYZdq{}}\PY{l+s+s2}{Jules Maigret}\PY{l+s+s2}{\PYZdq{}}
\end{Verbatim}


    \begin{Verbatim}[commandchars=\\\{\}]
{\color{incolor}In [{\color{incolor}11}]:} \PY{c+c1}{\PYZsh{} ainsi que l\PYZsq{}alphabet grec}
         \PY{k+kn}{from} \PY{n+nn}{math} \PY{k}{import} \PY{n}{cos}\PY{p}{,} \PY{n}{pi} \PY{k}{as} \PY{n}{π}
         \PY{n}{θ} \PY{o}{=} \PY{n}{π} \PY{o}{/} \PY{l+m+mi}{4}
         \PY{n}{cos}\PY{p}{(}\PY{n}{θ}\PY{p}{)}
\end{Verbatim}


\begin{Verbatim}[commandchars=\\\{\}]
{\color{outcolor}Out[{\color{outcolor}11}]:} 0.7071067811865476
\end{Verbatim}
            
    Tous les caractères Unicode ne sont pas permis - heureusement car cela
serait source de confusion. Nous citons dans les références les
documents qui précisent quels sont exactement les caractères autorisés.

    \begin{Verbatim}[commandchars=\\\{\}]
{\color{incolor}In [{\color{incolor}12}]:} \PY{c+c1}{\PYZsh{} ce caractère n\PYZsq{}est pas autorisé, car il}
         \PY{c+c1}{\PYZsh{} est considéré comme un signe mathématique (produit)}
         \PY{err}{∏} \PY{o}{=} \PY{l+m+mi}{10}
\end{Verbatim}


    \begin{Verbatim}[commandchars=\\\{\}]

          File "<ipython-input-12-4b0589d77bdc>", line 3
        ∏ = 10
        \^{}
    SyntaxError: invalid character in identifier


    \end{Verbatim}

    \begin{Verbatim}[commandchars=\\\{\}]
{\color{incolor}In [{\color{incolor}13}]:} \PY{c+c1}{\PYZsh{} ce caractère est encore différent, c\PYZsq{}est aussi}
         \PY{c+c1}{\PYZsh{} un pi grec mais pas le même, cette fois\PYZhy{}ci}
         \PY{c+c1}{\PYZsh{} c\PYZsq{}est un nom de variable acceptable mais }
         \PY{c+c1}{\PYZsh{} il n\PYZsq{}est pas défini}
         \PY{n}{∏}
\end{Verbatim}


    \begin{Verbatim}[commandchars=\\\{\}]

        ---------------------------------------------------------------------------

        NameError                                 Traceback (most recent call last)

        <ipython-input-13-42996045225e> in <module>()
          3 \# c'est un nom de variable acceptable mais
          4 \# il n'est pas défini
    ----> 5 Π
    

        NameError: name 'Π' is not defined

    \end{Verbatim}

    \hypertarget{conseil}{%
\paragraph{Conseil}\label{conseil}}

Il est \textbf{très vivement} recommandé~:

\begin{itemize}
\tightlist
\item
  tout d'abord de coder \textbf{en anglais}~;
\item
  ensuite de \textbf{ne pas} définir des identificateurs avec des
  caractères non ASCII, dans toute la mesure du possible~, voyez par
  exemple la confusion que peut créer le fait de nommer un
  identificateur π ou Π ou ∏~;
\item
  enfin si vous utilisez un encodage autre que UTF-8, vous
  \textbf{devez} bien \textbf{spécifier l'encodage} utilisé dans votre
  fichier source~; nous y reviendrons en deuxième semaine.
\end{itemize}

    \hypertarget{pour-en-savoir-plus}{%
\subsubsection{Pour en savoir plus}\label{pour-en-savoir-plus}}

    Pour les esprits curieux, Guido van Rossum, le fondateur de Python, est
le co-auteur d'un document qui décrit les conventions de codage à
utiliser dans la bibliothèque standard Python. Ces règles sont plus
restrictives que ce que le langage permet de faire, mais constituent une
lecture intéressante si vous projetez d'écrire beaucoup de Python.\\

Voir dans le PEP 008
\href{http://legacy.python.org/dev/peps/pep-0008/\#descriptive-naming-styles}{la
section consacrée aux règles de nommage - (en anglais)}\\

    Voir enfin, au sujet des caractères exotiques dans les identificateurs~:

\begin{itemize}
\tightlist
\item
  \href{https://www.python.org/dev/peps/pep-3131/}{https://www.python.org/dev/peps/pep-3131/} qui définit les caractères
  exotiques autorisés, et qui repose à son tour sur
\item
  \href{http://www.unicode.org/reports/tr31/}{http://www.unicode.org/reports/tr31/} (très technique~!)
\end{itemize}
        
    
    
    

    

    \hypertarget{les-mots-cluxe9s-de-python}{%
\section{Les mots-clés de Python}\label{les-mots-cluxe9s-de-python}}

    \hypertarget{mots-ruxe9servuxe9s}{%
\subsubsection{Mots réservés}\label{mots-ruxe9servuxe9s}}

    Il existe en Python certains mots spéciaux, qu'on appelle des mots-clés,
ou \emph{keywords} en anglais, qui sont réservés et \textbf{ne peuvent
pas être utilisés} comme identifiants, c'est-à-dire comme un nom de
variable.

    C'est le cas par exemple pour l'instruction \texttt{if}, que nous
verrons prochainement, qui permet bien entendu d'exécuter tel ou tel
code selon le résultat d'un test.

    \begin{Verbatim}[commandchars=\\\{\}]
{\color{incolor}In [{\color{incolor}1}]:} \PY{n}{variable} \PY{o}{=} \PY{l+m+mi}{15}
        \PY{k}{if} \PY{n}{variable} \PY{o}{\PYZlt{}}\PY{o}{=} \PY{l+m+mi}{10}\PY{p}{:}
            \PY{n+nb}{print}\PY{p}{(}\PY{l+s+s2}{\PYZdq{}}\PY{l+s+s2}{en dessous de la moyenne}\PY{l+s+s2}{\PYZdq{}}\PY{p}{)}
        \PY{k}{else}\PY{p}{:}
            \PY{n+nb}{print}\PY{p}{(}\PY{l+s+s2}{\PYZdq{}}\PY{l+s+s2}{au dessus}\PY{l+s+s2}{\PYZdq{}}\PY{p}{)}
\end{Verbatim}


    \begin{Verbatim}[commandchars=\\\{\}]
au dessus

    \end{Verbatim}

    À cause de la présence de cette instruction dans le langage, il n'est
pas autorisé d'appeler une variable \texttt{if}.

    \begin{Verbatim}[commandchars=\\\{\}]
{\color{incolor}In [{\color{incolor} }]:} \PY{c+c1}{\PYZsh{} NOTE}
        \PY{c+c1}{\PYZsh{} automatic execution has skipped execution of this cell}
        
        \PY{c+c1}{\PYZsh{} interdit, if est un mot\PYZhy{}clé}
        \PY{k}{if} \PY{o}{=} \PY{l+m+mi}{1}
\end{Verbatim}


    \hypertarget{liste-compluxe8te}{%
\subsubsection{Liste complète}\label{liste-compluxe8te}}

    Voici la liste complète des mots-clés~:

    \begin{longtable}[]{@{}lllll@{}}
\toprule
~ & ~ & ~ & ~ & ~\tabularnewline
\midrule
\endhead
\textbf{False} & await & else & import & pass\tabularnewline
\textbf{None} & break & except & in & raise\tabularnewline
\textbf{True} & class & finally & is & return\tabularnewline
and & continue & for & lambda & try\tabularnewline
as & def & from & \textbf{nonlocal} & while\tabularnewline
assert & del & global & not & with\tabularnewline
async & elif & if & or & yield\tabularnewline
\bottomrule
\end{longtable}

    Nous avons indiqué \textbf{en gras} les nouveautés \textbf{par rapport à
Python 2} (sachant que réciproquement \texttt{exec} et \texttt{print}
ont perdu leur statut de mot-clé depuis Python 2, ce sont maintenant des
fonctions).

    Il vous faudra donc y prêter attention, surtout au début, mais avec un
tout petit peu d'habitude vous saurez rapidement les éviter.

Vous remarquerez aussi que tous les bons éditeurs de texte supportant du
code Python vont colorer les mots-clés différemment des variables. Par
exemple, IDLE colorie les mots-clés en orange, vous pouvez donc très
facilement vous rendre compte que vous allez, par erreur, en utiliser un
comme nom de variable.

Cette fonctionnalité, dite de \emph{coloration syntaxique}, permet
d'identifier d'un coup d'œil, grâce à un code de couleur, le rôle des
différents éléments de votre code~: variables, mots-clés, etc. D'une
manière générale, nous vous déconseillons fortement d'utiliser un
éditeur de texte qui n'offre pas cette fonctionnalité de coloration
syntaxique.

    \hypertarget{pour-en-savoir-plus}{%
\subsubsection{Pour en savoir plus}\label{pour-en-savoir-plus}}

    On peut se reporter à cette page~:

https://docs.python.org/3/reference/lexical\_analysis.html\#keywords


    % Add a bibliography block to the postdoc
    
    
    

        
    
    
    

    

    \hypertarget{un-peu-de-calcul-sur-les-types}{%
\section{Un peu de calcul sur les
types}\label{un-peu-de-calcul-sur-les-types}}

    \hypertarget{compluxe9ment---niveau-basique}{%
\subsection{Complément - niveau
basique}\label{compluxe9ment---niveau-basique}}

    \hypertarget{la-fonction-type}{%
\subsubsection{\texorpdfstring{La fonction
\texttt{type}}{La fonction type}}\label{la-fonction-type}}

    Nous avons vu dans la vidéo que chaque objet possède un type. On peut
très simplement accéder au type d'un objet en appelant une fonction
\emph{built-in}, c'est-à-dire prédéfinie dans Python, qui s'appelle, eh
bien oui, \texttt{type}.

    On l'utilise tout simplement comme ceci~:

    \begin{Verbatim}[commandchars=\\\{\}]
{\color{incolor}In [{\color{incolor}1}]:} \PY{n+nb}{type}\PY{p}{(}\PY{l+m+mi}{1}\PY{p}{)}
\end{Verbatim}


\begin{Verbatim}[commandchars=\\\{\}]
{\color{outcolor}Out[{\color{outcolor}1}]:} int
\end{Verbatim}
            
    \begin{Verbatim}[commandchars=\\\{\}]
{\color{incolor}In [{\color{incolor}2}]:} \PY{n+nb}{type}\PY{p}{(}\PY{l+s+s1}{\PYZsq{}}\PY{l+s+s1}{spam}\PY{l+s+s1}{\PYZsq{}}\PY{p}{)}
\end{Verbatim}


\begin{Verbatim}[commandchars=\\\{\}]
{\color{outcolor}Out[{\color{outcolor}2}]:} str
\end{Verbatim}
            
    Cette fonction est assez peu utilisée par les programmeurs expérimentés,
mais va nous être utile à bien comprendre le langage, notamment pour
manipuler les valeurs numériques.

    \hypertarget{types-variables-et-objets}{%
\subsubsection{Types, variables et
objets}\label{types-variables-et-objets}}

    On a vu également que le type est attaché \textbf{à l'objet} et non à la
variable.

    \begin{Verbatim}[commandchars=\\\{\}]
{\color{incolor}In [{\color{incolor}3}]:} \PY{n}{x} \PY{o}{=} \PY{l+m+mi}{1}
        \PY{n+nb}{type}\PY{p}{(}\PY{n}{x}\PY{p}{)}
\end{Verbatim}


\begin{Verbatim}[commandchars=\\\{\}]
{\color{outcolor}Out[{\color{outcolor}3}]:} int
\end{Verbatim}
            
    \begin{Verbatim}[commandchars=\\\{\}]
{\color{incolor}In [{\color{incolor}4}]:} \PY{c+c1}{\PYZsh{} la variable x peut référencer un objet de n\PYZsq{}importe quel type}
        
        \PY{n}{x} \PY{o}{=} \PY{p}{[}\PY{l+m+mi}{1}\PY{p}{,} \PY{l+m+mi}{2}\PY{p}{,} \PY{l+m+mi}{3}\PY{p}{]}
        \PY{n+nb}{type}\PY{p}{(}\PY{n}{x}\PY{p}{)}
\end{Verbatim}


\begin{Verbatim}[commandchars=\\\{\}]
{\color{outcolor}Out[{\color{outcolor}4}]:} list
\end{Verbatim}
            
    \hypertarget{compluxe9ment---niveau-avancuxe9}{%
\subsection{Complément - niveau
avancé}\label{compluxe9ment---niveau-avancuxe9}}

    \hypertarget{la-fonction-isinstance}{%
\subsubsection{\texorpdfstring{La fonction
\texttt{isinstance}}{La fonction isinstance}}\label{la-fonction-isinstance}}

    Une autre fonction prédéfinie, voisine de \texttt{type} mais plus utile
dans la pratique, est la fonction \texttt{isinstance} qui permet de
savoir si un objet est d'un type donné. Par exemple~:

    \begin{Verbatim}[commandchars=\\\{\}]
{\color{incolor}In [{\color{incolor}5}]:} \PY{n+nb}{isinstance}\PY{p}{(}\PY{l+m+mi}{23}\PY{p}{,} \PY{n+nb}{int}\PY{p}{)}
\end{Verbatim}


\begin{Verbatim}[commandchars=\\\{\}]
{\color{outcolor}Out[{\color{outcolor}5}]:} True
\end{Verbatim}
            
    À la vue de ce seul exemple, on pourrait penser que \texttt{isinstance}
est presque identique à \texttt{type}~; en réalité elle est un peu plus
élaborée, notamment pour la programmation objet et l'héritage, nous
aurons l'occasion d'y revenir.

    On remarque ici en passant que la variable \texttt{int} est connue de
Python alors que nous ne l'avons pas définie. Il s'agit d'une variable
prédéfinie, qui désigne le type des entiers, que nous étudierons très
bientôt.

    Pour conclure sur \texttt{isinstance}, cette fonction est utile en
pratique précisément parce que Python est à typage dynamique. Aussi il
est souvent utile de s'assurer qu'une variable passée à une fonction est
du (ou des) type(s) attendu(s), puisque contrairement à un langage typé
statiquement comme C++, on n'a aucune garantie de ce genre à
l'exécution. À nouveau, nous aurons l'occasion de revenir sur ce point.


    % Add a bibliography block to the postdoc
    
    
    

        \hypertarget{gestion-de-la-muxe9moire}{%
\section{Gestion de la mémoire}\label{gestion-de-la-muxe9moire}}

    \hypertarget{compluxe9ment---niveau-basique}{%
\subsection{Complément - niveau
basique}\label{compluxe9ment---niveau-basique}}

    L'objet de ce complément est de vous montrer qu'avec Python vous n'avez
pas à vous préoccuper de la mémoire. Pour expliquer la notion de gestion
de la mémoire, il nous faut donner un certain nombre de détails sur
d'autres langages comme C et C++. Si vous souhaitez suivre ce cours à un
niveau basique vous pouvez ignorer ce complément et seulement retenir
que Python se charge de tout pour vous~:)

    \hypertarget{compluxe9ment---niveau-intermuxe9diaire}{%
\subsection{Complément - niveau
intermédiaire}\label{compluxe9ment---niveau-intermuxe9diaire}}

    \hypertarget{langages-de-bas-niveau}{%
\subsubsection{Langages de bas niveau}\label{langages-de-bas-niveau}}

    Dans un langage traditionnel de bas niveau comme C ou C++, le
programmeur est en charge de l'allocation - et donc de la libération -
de la mémoire.\\

Ce qui signifie que, sauf pour les valeurs stockées dans la pile, le
programmeur est amené~:
\begin{itemize}
	\item 
	à réclamer de la mémoire au système
	d'exploitation en appelant explicitement \texttt{malloc} (C) ou
	\texttt{new} (C++)~;
	\item
	et réciproquement à rendre cette mémoire au
	système d'exploitation lorsqu'elle n'est plus utilisée, en appelant
	\texttt{free} (C) ou \texttt{delete} (C++).
\end{itemize}

    Avec ce genre de langage, la gestion de la mémoire est un aspect
important de la programmation. Ce modèle offre une grande flexibilité,
mais au prix d'un coût élevé en matière de vitesse de développement.\\

En effet, il est assez facile d'oublier de libérer la mémoire après
usage, ce qui peut conduire à épuiser les ressources disponibles. À
l'inverse, utiliser une zone mémoire non allouée peut conduire à des
bugs très difficiles à localiser et à des problèmes de sécurité majeurs.
Notons qu'une grande partie des attaques en informatique reposent sur
l'exploitation d'erreurs de gestion de la mémoire.

    \hypertarget{langages-de-haut-niveau}{%
\subsubsection{Langages de haut niveau}\label{langages-de-haut-niveau}}

    Pour toutes ces raisons, avec un langage de plus haut niveau comme
Python, le programmeur est libéré de cet aspect de la programmation.\\

Pour anticiper un peu sur le cours des semaines suivantes, voici ce que
vous pouvez garder en tête s'agissant de la gestion mémoire en Python~:
* vous créez vos objets au fur et à mesure de vos besoins~;

\begin{itemize}
\item
  vous n'avez pas besoin de les libérer explicitement, le
  ``\emph{Garbage Collector}'' de Python va s'en charger pour recycler
  la mémoire lorsque c'est possible~;
\item
  Python a tendance à être assez gourmand en mémoire, comparé à un
  langage de bas niveau, car tout est objet et chaque objet est assorti
  de \emph{méta-informations} qui occupent une place non négligeable.
  Par exemple, chaque objet possède au minimum~:

  \begin{itemize}
  \tightlist
  \item
    une référence vers son type - c'est le prix du typage dynamique~;
  \item
    un compteur de références - le nombre d'autres valeurs (variables ou
    objets) qui pointent vers l'objet, cette information est notamment
    utilisée, précisément, par le \emph{Garbage Collector} pour
    déterminer si la mémoire utilisée par un objet peut être libérée ou
    non.
  \end{itemize}
\item
  un certain nombre de types prédéfinis et non mutables sont implémentés
  en Python comme des \emph{singletons}, c'est-à-dire qu'un seul objet
  est créé et partagé, c'est le cas par exemple pour les petits entiers
  et les chaînes de caractères, on en reparlera~;
\item
  lorsqu'on implémente une classe, il est possible de lui conférer cette
  caractéristique de singleton, de manière à optimiser la mémoire
  nécessaire pour exécuter un programme.
\end{itemize}
        \hypertarget{typages-statique-et-dynamique}{%
\section{Typages statique et
dynamique}\label{typages-statique-et-dynamique}}

    \hypertarget{compluxe9ment---niveau-intermuxe9diaire}{%
\subsection{Complément - niveau
intermédiaire}\label{compluxe9ment---niveau-intermuxe9diaire}}

    Parmi les langages typés, on distingue les langages à typage statique et
ceux à typage dynamique. Ce notebook tente d'éclaircir ces notions pour
ceux qui n'y sont pas familiers.

    \hypertarget{typage-statique}{%
\subsubsection{Typage statique}\label{typage-statique}}

    À une extrémité du spectre, on trouve les langages compilés, dits à
typage statique, comme par exemple C ou C++.

En C on écrira, par exemple, une version simpliste de la fonction
factoriel comme ceci~:

    \begin{Shaded}
\begin{Highlighting}[]
\DataTypeTok{int}\NormalTok{ factoriel(}\DataTypeTok{int}\NormalTok{ n) \{}
    \DataTypeTok{int}\NormalTok{ result = }\DecValTok{1}\NormalTok{;}
    \ControlFlowTok{for}\NormalTok{ (}\DataTypeTok{int}\NormalTok{ loop = }\DecValTok{1}\NormalTok{; loop <= n; loop++)}
\NormalTok{        result *= loop;}
    \ControlFlowTok{return}\NormalTok{ result;}
\NormalTok{\}}
\end{Highlighting}
\end{Shaded}

    Comme vous pouvez le voir - ou le deviner - toutes les
\textbf{variables} utilisées ici (comme par exemple \texttt{n},
\texttt{result} et \texttt{loop}) sont typées~:

\begin{itemize}
\tightlist
\item
  on doit appeler \texttt{factoriel} avec un argument \texttt{n} qui
  doit être un entier (\texttt{int} est le nom du type entier)~;
\item
  les variables internes \texttt{result} et \texttt{loop} sont de type
  entier~;
\item
  \texttt{factoriel} retourne une valeur de type entier.
\end{itemize}

    Ces informations de type ont essentiellement trois fonctions~:

\begin{itemize}
\tightlist
\item
  en premier lieu, elles sont nécessaires au compilateur. En C si le
  programmeur ne précisait pas que \texttt{result} est de type entier,
  le compilateur n'aurait pas suffisamment d'éléments pour générer le
  code assembleur correspondant~;
\item
  en contrepartie, le programmeur a un contrôle très fin de l'usage
  qu'il fait de la mémoire et du matériel. Il peut choisir d'utiliser un
  entier sur 32 ou 64 bits, signé ou pas, ou construire avec
  \texttt{struct} et \texttt{union} un arrangement de ses données~;
\item
  enfin, et surtout, ces informations de type permettent de faire un
  contrôle \emph{a priori} de la validité du programme, par exemple, si
  à un autre endroit dans le code on trouve~:
\end{itemize}

    \begin{Shaded}
\begin{Highlighting}[]
\PreprocessorTok{#include }\ImportTok{<stdio.h>}

\DataTypeTok{int}\NormalTok{ main(}\DataTypeTok{int}\NormalTok{ argc, }\DataTypeTok{char}\NormalTok{ *argv[]) \{}
    \CommentTok{/* le premier argument de la ligne de commande est argv[1] */}
    \DataTypeTok{char}\NormalTok{ *input = argv[}\DecValTok{1}\NormalTok{];}
    \CommentTok{/* calculer son factoriel et afficher le résultat */}
\NormalTok{    printf(}\StringTok{"Factoriel (%s) = %d}\SpecialCharTok{\textbackslash{}n}\StringTok{"}\NormalTok{, input, factoriel(input));}
    \CommentTok{/*                                               ^^^^^}
\CommentTok{     * ici on appelle factoriel avec une entrée de type 'chaîne de caractères' */}
\NormalTok{\}}
\end{Highlighting}
\end{Shaded}

    alors le compilateur va remarquer qu'on essaie d'appeler
\texttt{factoriel} avec comme argument \texttt{input} qui, pour faire
simple, est une chaîne de caractères et comme \texttt{factoriel}
s'attend à recevoir un entier, ce programme n'a aucune chance de
compiler.\\

On parle alors de \textbf{typage statique}, en ce sens que chaque
\textbf{variable} a exactement un type qui est défini par le programmeur
une bonne fois pour toutes.\\

    C'est ce qu'on appelle le \textbf{contrôle de type}, ou
\emph{type-checking} en anglais. Si on ignore le point sur le contrôle
fin de la mémoire, qui n'est pas crucial à notre sujet, ce modèle de
contrôle de type présente~:

\begin{itemize}
\tightlist
\item
  l'\textbf{inconvénient} de demander davantage au programmeur (je fais
  abstraction, à ce stade et pour simplifier, de
  \href{https://en.wikipedia.org/wiki/Type_inference}{langages à
  inférence de types} comme ML et Haskell)~;
\item
  et l'\textbf{avantage} de permettre un contrôle étendu, et surtout
  précoce (avant même de l'exécuter), de la bonne correction du
  programme.
\end{itemize}

    Cela étant dit, le typage statique en C n'empêche pas le programmeur
débutant d'essayer d'écrire dans la mémoire à partir d'un pointeur
\texttt{NULL} - et le programme de s'interrompre brutalement. Il faut
être conscient des limites du typage statique.

    \hypertarget{typage-dynamique}{%
\subsubsection{Typage dynamique}\label{typage-dynamique}}

    À l'autre bout du spectre, on trouve des langages comme, eh bien,
Python.\\

    Pour comprendre cette notion de typage dynamique, regardons la fonction
suivante \texttt{somme}.

    \begin{Verbatim}[commandchars=\\\{\}]
{\color{incolor}In [{\color{incolor}1}]:} \PY{k}{def} \PY{n+nf}{somme}\PY{p}{(}\PY{o}{*}\PY{n}{largs}\PY{p}{)}\PY{p}{:}
            \PY{l+s+s2}{\PYZdq{}}\PY{l+s+s2}{retourne la somme de tous ses arguments}\PY{l+s+s2}{\PYZdq{}}
            \PY{k}{if} \PY{o+ow}{not} \PY{n}{largs}\PY{p}{:}
                \PY{k}{return} \PY{l+m+mi}{0}
            \PY{n}{result} \PY{o}{=} \PY{n}{largs}\PY{p}{[}\PY{l+m+mi}{0}\PY{p}{]}
            \PY{k}{for} \PY{n}{i} \PY{o+ow}{in} \PY{n+nb}{range}\PY{p}{(}\PY{l+m+mi}{1}\PY{p}{,} \PY{n+nb}{len}\PY{p}{(}\PY{n}{largs}\PY{p}{)}\PY{p}{)}\PY{p}{:}
                \PY{n}{result} \PY{o}{+}\PY{o}{=} \PY{n}{largs}\PY{p}{[}\PY{n}{i}\PY{p}{]}
            \PY{k}{return} \PY{n}{result}
\end{Verbatim}


    Naturellement, vous n'êtes pas à ce stade en mesure de comprendre le
fonctionnement intime de la fonction. Mais vous pouvez tout de même
l'utiliser~:

    \begin{Verbatim}[commandchars=\\\{\}]
{\color{incolor}In [{\color{incolor}2}]:} \PY{n}{somme}\PY{p}{(}\PY{l+m+mi}{12}\PY{p}{,} \PY{l+m+mi}{14}\PY{p}{,} \PY{l+m+mi}{300}\PY{p}{)}
\end{Verbatim}


\begin{Verbatim}[commandchars=\\\{\}]
{\color{outcolor}Out[{\color{outcolor}2}]:} 326
\end{Verbatim}
            
    \begin{Verbatim}[commandchars=\\\{\}]
{\color{incolor}In [{\color{incolor}3}]:} \PY{n}{liste1} \PY{o}{=} \PY{p}{[}\PY{l+s+s1}{\PYZsq{}}\PY{l+s+s1}{a}\PY{l+s+s1}{\PYZsq{}}\PY{p}{,} \PY{l+s+s1}{\PYZsq{}}\PY{l+s+s1}{b}\PY{l+s+s1}{\PYZsq{}}\PY{p}{,} \PY{l+s+s1}{\PYZsq{}}\PY{l+s+s1}{c}\PY{l+s+s1}{\PYZsq{}}\PY{p}{]}
        \PY{n}{liste2} \PY{o}{=} \PY{p}{[}\PY{l+m+mi}{0}\PY{p}{,} \PY{l+m+mi}{20}\PY{p}{,} \PY{l+m+mi}{30}\PY{p}{]}
        \PY{n}{liste3} \PY{o}{=} \PY{p}{[}\PY{l+s+s1}{\PYZsq{}}\PY{l+s+s1}{spam}\PY{l+s+s1}{\PYZsq{}}\PY{p}{,} \PY{l+s+s1}{\PYZsq{}}\PY{l+s+s1}{eggs}\PY{l+s+s1}{\PYZsq{}}\PY{p}{]}
        \PY{n}{somme}\PY{p}{(}\PY{n}{liste1}\PY{p}{,} \PY{n}{liste2}\PY{p}{,} \PY{n}{liste3}\PY{p}{)}
\end{Verbatim}


\begin{Verbatim}[commandchars=\\\{\}]
{\color{outcolor}Out[{\color{outcolor}3}]:} ['a', 'b', 'c', 0, 20, 30, 'spam', 'eggs']
\end{Verbatim}
            
    Vous pouvez donc constater que \texttt{somme} peut fonctionner avec des
objets de types différents. En fait, telle qu'elle est écrite, elle va
fonctionner s'il est possible de faire \texttt{+} entre ses arguments.
Ainsi, par exemple, on pourrait même faire~:

    \begin{Verbatim}[commandchars=\\\{\}]
{\color{incolor}In [{\color{incolor}4}]:} \PY{c+c1}{\PYZsh{} Python sait faire + entre deux chaînes de caractères}
        \PY{n}{somme}\PY{p}{(}\PY{l+s+s1}{\PYZsq{}}\PY{l+s+s1}{abc}\PY{l+s+s1}{\PYZsq{}}\PY{p}{,} \PY{l+s+s1}{\PYZsq{}}\PY{l+s+s1}{def}\PY{l+s+s1}{\PYZsq{}}\PY{p}{)}
\end{Verbatim}


\begin{Verbatim}[commandchars=\\\{\}]
{\color{outcolor}Out[{\color{outcolor}4}]:} 'abcdef'
\end{Verbatim}
            
    Mais par contre on ne pourrait pas faire

    \begin{Verbatim}[commandchars=\\\{\}]
{\color{incolor}In [{\color{incolor}5}]:} \PY{c+c1}{\PYZsh{} ceci va déclencher une exception à l\PYZsq{}exécution}
        \PY{n}{somme}\PY{p}{(}\PY{l+m+mi}{12}\PY{p}{,} \PY{p}{[}\PY{l+m+mi}{1}\PY{p}{,} \PY{l+m+mi}{2}\PY{p}{,} \PY{l+m+mi}{3}\PY{p}{]}\PY{p}{)}
\end{Verbatim}


    \begin{Verbatim}[commandchars=\\\{\}]

        ---------------------------------------------------------------------------

        TypeError                                 Traceback (most recent call last)

        <ipython-input-5-1b5269e9e129> in <module>()
          1 \# ceci va déclencher une exception à l'exécution
    ----> 2 somme(12, [1, 2, 3])
    

        <ipython-input-1-29005c25d5bb> in somme(*largs)
          5     result = largs[0]
          6     for i in range(1, len(largs)):
    ----> 7         result += largs[i]
          8     return result


        TypeError: unsupported operand type(s) for +=: 'int' and 'list'

    \end{Verbatim}

    Il est utile de remarquer que le typage de Python, qui existe bel et
bien comme on le verra, est qualifié de dynamique parce que le type est
attaché \textbf{à un objet} et non à la variable qui le référence. On
aura bien entendu l'occasion d'approfondir tout ça dans le cours.\\

    En Python, on fait souvent référence au typage sous l'appellation
\emph{duck typing}, de manière imagée~:\\

\begin{quote}
If it looks like a duck and quacks like a duck, it's a duck.\\
\end{quote}

    On voit qu'on se trouve dans une situation très différente de celle du
programmeur C/C++, en ce sens que~:

\begin{itemize}
\tightlist
\item
  à l'écriture du programme, il n'y aucun des surcoûts qu'on trouve avec
  C ou C++ en matière de définition de type~;
\item
  aucun contrôle de type n'est effectué \emph{a priori} par le langage
  au moment de la définition de la fonction \texttt{somme}~;
\item
  par contre au moment de l'exécution, s'il s'avère qu'on tente de faire
  une somme entre deux types qui ne peuvent pas être additionnés, comme
  ci-dessus avec un entier et une liste, le programme ne pourra pas se
  dérouler correctement.
\end{itemize}

    Il y a deux points de vue vis-à-vis de la question du typage.\\

Les gens habitués au \emph{typage statique} se plaignent du typage
dynamique en disant qu'on peut écrire des programmes faux et qu'on s'en
rend compte trop tard - à l'exécution.\\

À l'inverse les gens habitués au \emph{typage dynamique} font valoir que
le typage statique est très partiel, par exemple, en C si on essaie
d'écrire dans un pointeur \texttt{NULL}, le système d'exploitation ne le
permet pas et le programme sort tout aussi brutalement.\\

    Bref, selon le point de vue, le typage dynamique est vécu comme un
inconvénient (pas assez de bonnes propriétés détectées par le langage)
ou comme un avantage (pas besoin de passer du temps à déclarer le type
des variables, ni à faire des conversions pour satisfaire le
compilateur).\\

Vous remarquerez cependant à l'usage, qu'en matière de vitesse de
développement, les inconvénients du typage dynamique sont très largement
compensés par ses avantages.

    \hypertarget{type-hints}{%
\subsubsection{\texorpdfstring{\emph{Type
hints}}{Type hints}}\label{type-hints}}

    Signalons enfin que depuis python-3.5, il est \textbf{possible}
d'ajouter des annotations de type, pour expliciter les suppositions qui
sont faites par le programmeur pour le bon fonctionnement du code.\\

Nous aurons là encore l'occasion de détailler ce point dans le cours,
signalons simplement que ces annotations sont totalement optionnelles,
et que même lorsqu'elles sont présentes elles ne sont pas utilisées à
l'exécution par l'interpréteur. L'idée est plutôt de permettre à des
outils externes, \href{http://www.mypy-lang.org}{comme par exemple
\texttt{mypy}}, d'effectuer des contrôles plus poussés concernant la
correction du programme.
        
    
    
    

    

    \hypertarget{utiliser-python-comme-une-calculette}{%
\section{Utiliser Python comme une
calculette}\label{utiliser-python-comme-une-calculette}}

    Lorsque vous démarrez l'interprète Python, vous disposez en fait d'une
calculette, par exemple, vous pouvez taper~:

    \begin{Verbatim}[commandchars=\\\{\}]
{\color{incolor}In [{\color{incolor}1}]:} \PY{l+m+mi}{20} \PY{o}{*} \PY{l+m+mi}{60}
\end{Verbatim}


\begin{Verbatim}[commandchars=\\\{\}]
{\color{outcolor}Out[{\color{outcolor}1}]:} 1200
\end{Verbatim}
            
    Les règles de \textbf{priorité} entre les opérateurs sont habituelles,
les produits et divisions sont évalués en premier, ensuite les sommes et
soustractions~:

    \begin{Verbatim}[commandchars=\\\{\}]
{\color{incolor}In [{\color{incolor}2}]:} \PY{l+m+mi}{2} \PY{o}{*} \PY{l+m+mi}{30} \PY{o}{+} \PY{l+m+mi}{10} \PY{o}{*} \PY{l+m+mi}{5}
\end{Verbatim}


\begin{Verbatim}[commandchars=\\\{\}]
{\color{outcolor}Out[{\color{outcolor}2}]:} 110
\end{Verbatim}
            
    De manière générale, il est recommandé de bien parenthéser ses
expressions. De plus, les parenthèses facilitent la lecture
d'expressions complexes.

Par exemple, il vaut mieux écrire ce qui suit, qui est équivalent mais
plus lisible~:

    \begin{Verbatim}[commandchars=\\\{\}]
{\color{incolor}In [{\color{incolor}3}]:} \PY{p}{(}\PY{l+m+mi}{2} \PY{o}{*} \PY{l+m+mi}{30}\PY{p}{)} \PY{o}{+} \PY{p}{(}\PY{l+m+mi}{10} \PY{o}{*} \PY{l+m+mi}{5}\PY{p}{)}
\end{Verbatim}


\begin{Verbatim}[commandchars=\\\{\}]
{\color{outcolor}Out[{\color{outcolor}3}]:} 110
\end{Verbatim}
            
    Attention, en Python la division \texttt{/} est une division naturelle~:

    \begin{Verbatim}[commandchars=\\\{\}]
{\color{incolor}In [{\color{incolor}4}]:} \PY{l+m+mi}{48} \PY{o}{/} \PY{l+m+mi}{5}
\end{Verbatim}


\begin{Verbatim}[commandchars=\\\{\}]
{\color{outcolor}Out[{\color{outcolor}4}]:} 9.6
\end{Verbatim}
            
    Rappelez-vous des opérateurs suivants qui sont très pratiques~:

\begin{longtable}[]{@{}ll@{}}
\toprule
code & opération\tabularnewline
\midrule
\endhead
\texttt{//} & quotient\tabularnewline
\texttt{\%} & modulo\tabularnewline
\texttt{**} & puissance\tabularnewline
\bottomrule
\end{longtable}

    \begin{Verbatim}[commandchars=\\\{\}]
{\color{incolor}In [{\color{incolor}5}]:} \PY{c+c1}{\PYZsh{} calculer un quotient}
        \PY{l+m+mi}{48} \PY{o}{/}\PY{o}{/} \PY{l+m+mi}{5}
\end{Verbatim}


\begin{Verbatim}[commandchars=\\\{\}]
{\color{outcolor}Out[{\color{outcolor}5}]:} 9
\end{Verbatim}
            
    \begin{Verbatim}[commandchars=\\\{\}]
{\color{incolor}In [{\color{incolor}6}]:} \PY{c+c1}{\PYZsh{} modulo (le reste de la division par)}
        \PY{l+m+mi}{48} \PY{o}{\PYZpc{}} \PY{l+m+mi}{5}
\end{Verbatim}


\begin{Verbatim}[commandchars=\\\{\}]
{\color{outcolor}Out[{\color{outcolor}6}]:} 3
\end{Verbatim}
            
    \begin{Verbatim}[commandchars=\\\{\}]
{\color{incolor}In [{\color{incolor}7}]:} \PY{c+c1}{\PYZsh{} puissance}
        \PY{l+m+mi}{2} \PY{o}{*}\PY{o}{*} \PY{l+m+mi}{10}
\end{Verbatim}


\begin{Verbatim}[commandchars=\\\{\}]
{\color{outcolor}Out[{\color{outcolor}7}]:} 1024
\end{Verbatim}
            
    Vous pouvez facilement faire aussi des calculs sur les complexes.
Souvenez-vous seulement que la constante complexe que nous notons
\texttt{i} en français se note \texttt{j} en Python, ce choix a été fait
par \href{https://fr.wikipedia.org/wiki/Benevolent_Dictator_for_Life}{le
BDFL} - alias Guido van Rossum - pour des raisons de lisibilité~:

    \begin{Verbatim}[commandchars=\\\{\}]
{\color{incolor}In [{\color{incolor}8}]:} \PY{c+c1}{\PYZsh{} multiplication de deux nombres complexes}
        \PY{p}{(}\PY{l+m+mi}{2} \PY{o}{+} \PY{l+m+mi}{3}\PY{n}{j}\PY{p}{)} \PY{o}{*} \PY{l+m+mf}{2.5}\PY{n}{j}
\end{Verbatim}


\begin{Verbatim}[commandchars=\\\{\}]
{\color{outcolor}Out[{\color{outcolor}8}]:} (-7.5+5j)
\end{Verbatim}
            
    Aussi, pour entrer ce nombre complexe \texttt{j}, il faut toujours le
faire précéder d'un nombre, donc ne pas entrer simplement \texttt{j}
(qui serait compris comme un nom de variable, nous allons voir ça tout
de suite) mais plutôt \texttt{1j} ou encore \texttt{1.j}, comme ceci~:

    \begin{Verbatim}[commandchars=\\\{\}]
{\color{incolor}In [{\color{incolor}9}]:} \PY{l+m+mi}{1}\PY{n}{j} \PY{o}{*} \PY{l+m+mf}{1.}\PY{n}{j}
\end{Verbatim}


\begin{Verbatim}[commandchars=\\\{\}]
{\color{outcolor}Out[{\color{outcolor}9}]:} (-1+0j)
\end{Verbatim}
            
    \hypertarget{utiliser-des-variables}{%
\subsubsection{Utiliser des variables}\label{utiliser-des-variables}}

    Il peut être utile de stocker un résultat qui sera utilisé plus tard, ou
de définir une valeur constante. Pour cela on utilise tout simplement
une affectation comme ceci~:

    \begin{Verbatim}[commandchars=\\\{\}]
{\color{incolor}In [{\color{incolor}10}]:} \PY{c+c1}{\PYZsh{} pour définir une variable il suffit de lui assigner une valeur}
         \PY{n}{largeur} \PY{o}{=} \PY{l+m+mi}{5}
\end{Verbatim}


    \begin{Verbatim}[commandchars=\\\{\}]
{\color{incolor}In [{\color{incolor}11}]:} \PY{c+c1}{\PYZsh{} une fois la variable définie, on peut l\PYZsq{}utiliser, ici comme un nombre}
         \PY{n}{largeur} \PY{o}{*} \PY{l+m+mi}{20}
\end{Verbatim}


\begin{Verbatim}[commandchars=\\\{\}]
{\color{outcolor}Out[{\color{outcolor}11}]:} 100
\end{Verbatim}
            
    \begin{Verbatim}[commandchars=\\\{\}]
{\color{incolor}In [{\color{incolor}12}]:} \PY{c+c1}{\PYZsh{} après quoi bien sûr la variable reste inchangée}
         \PY{n}{largeur} \PY{o}{*} \PY{l+m+mi}{10}
\end{Verbatim}


\begin{Verbatim}[commandchars=\\\{\}]
{\color{outcolor}Out[{\color{outcolor}12}]:} 50
\end{Verbatim}
            
    Pour les symboles mathématiques, on peut utiliser la même technique~:

    \begin{Verbatim}[commandchars=\\\{\}]
{\color{incolor}In [{\color{incolor}13}]:} \PY{c+c1}{\PYZsh{} pour définir un réel, on utilise le point au lieu d\PYZsq{}une virgule en français}
         \PY{n}{pi} \PY{o}{=} \PY{l+m+mf}{3.14159}
         \PY{l+m+mi}{2} \PY{o}{*} \PY{n}{pi} \PY{o}{*} \PY{l+m+mi}{10}
\end{Verbatim}


\begin{Verbatim}[commandchars=\\\{\}]
{\color{outcolor}Out[{\color{outcolor}13}]:} 62.8318
\end{Verbatim}
            
    Pour les valeurs spéciales comme \(\pi\), on peut utiliser les valeurs
prédéfinies par la bibliothèque mathématique de Python. En anticipant un
peu sur la notion d'importation que nous approfondirons plus tard, on
peut écrire~:

    \begin{Verbatim}[commandchars=\\\{\}]
{\color{incolor}In [{\color{incolor}14}]:} \PY{k+kn}{from} \PY{n+nn}{math} \PY{k}{import} \PY{n}{e}\PY{p}{,} \PY{n}{pi}
\end{Verbatim}


    Et ainsi imprimer les racines troisièmes de l'unité par la formule~:

\(r_n = e^{2i\pi \frac{n}{3}},\) pour \(n\in \{0,1,2\}\)

    \begin{Verbatim}[commandchars=\\\{\}]
{\color{incolor}In [{\color{incolor}15}]:} \PY{n}{n} \PY{o}{=} \PY{l+m+mi}{0}
         \PY{n+nb}{print}\PY{p}{(}\PY{l+s+s2}{\PYZdq{}}\PY{l+s+s2}{n=}\PY{l+s+s2}{\PYZdq{}}\PY{p}{,} \PY{n}{n}\PY{p}{,} \PY{l+s+s2}{\PYZdq{}}\PY{l+s+s2}{racine = }\PY{l+s+s2}{\PYZdq{}}\PY{p}{,} \PY{n}{e}\PY{o}{*}\PY{o}{*}\PY{p}{(}\PY{p}{(}\PY{l+m+mf}{2.}\PY{n}{j}\PY{o}{*}\PY{n}{pi}\PY{o}{*}\PY{n}{n}\PY{p}{)}\PY{o}{/}\PY{l+m+mi}{3}\PY{p}{)}\PY{p}{)}
         \PY{n}{n} \PY{o}{=} \PY{l+m+mi}{1}
         \PY{n+nb}{print}\PY{p}{(}\PY{l+s+s2}{\PYZdq{}}\PY{l+s+s2}{n=}\PY{l+s+s2}{\PYZdq{}}\PY{p}{,} \PY{n}{n}\PY{p}{,} \PY{l+s+s2}{\PYZdq{}}\PY{l+s+s2}{racine = }\PY{l+s+s2}{\PYZdq{}}\PY{p}{,} \PY{n}{e}\PY{o}{*}\PY{o}{*}\PY{p}{(}\PY{p}{(}\PY{l+m+mf}{2.}\PY{n}{j}\PY{o}{*}\PY{n}{pi}\PY{o}{*}\PY{n}{n}\PY{p}{)}\PY{o}{/}\PY{l+m+mi}{3}\PY{p}{)}\PY{p}{)}
         \PY{n}{n} \PY{o}{=} \PY{l+m+mi}{2}
         \PY{n+nb}{print}\PY{p}{(}\PY{l+s+s2}{\PYZdq{}}\PY{l+s+s2}{n=}\PY{l+s+s2}{\PYZdq{}}\PY{p}{,} \PY{n}{n}\PY{p}{,} \PY{l+s+s2}{\PYZdq{}}\PY{l+s+s2}{racine = }\PY{l+s+s2}{\PYZdq{}}\PY{p}{,} \PY{n}{e}\PY{o}{*}\PY{o}{*}\PY{p}{(}\PY{p}{(}\PY{l+m+mf}{2.}\PY{n}{j}\PY{o}{*}\PY{n}{pi}\PY{o}{*}\PY{n}{n}\PY{p}{)}\PY{o}{/}\PY{l+m+mi}{3}\PY{p}{)}\PY{p}{)}
\end{Verbatim}


    \begin{Verbatim}[commandchars=\\\{\}]
n= 0 racine =  (1+0j)
n= 1 racine =  (-0.4999999999999998+0.8660254037844388j)
n= 2 racine =  (-0.5000000000000004-0.8660254037844384j)

    \end{Verbatim}

    \textbf{Remarque~:} bien entendu il sera possible de faire ceci plus
simplement lorsque nous aurons vu les boucles \texttt{for}.

    \hypertarget{les-types}{%
\subsubsection{Les types}\label{les-types}}

    Ce qui change par rapport à une calculatrice standard est le fait que
les valeurs sont typées. Pour illustrer les trois types de nombres que
nous avons vus jusqu'ici~:

    \begin{Verbatim}[commandchars=\\\{\}]
{\color{incolor}In [{\color{incolor}16}]:} \PY{c+c1}{\PYZsh{} le type entier s\PYZsq{}appelle \PYZsq{}int\PYZsq{}}
         \PY{n+nb}{type}\PY{p}{(}\PY{l+m+mi}{3}\PY{p}{)}
\end{Verbatim}


\begin{Verbatim}[commandchars=\\\{\}]
{\color{outcolor}Out[{\color{outcolor}16}]:} int
\end{Verbatim}
            
    \begin{Verbatim}[commandchars=\\\{\}]
{\color{incolor}In [{\color{incolor}17}]:} \PY{c+c1}{\PYZsh{} le type flottant s\PYZsq{}appelle \PYZsq{}float\PYZsq{}}
         \PY{n+nb}{type}\PY{p}{(}\PY{l+m+mf}{3.5}\PY{p}{)}
\end{Verbatim}


\begin{Verbatim}[commandchars=\\\{\}]
{\color{outcolor}Out[{\color{outcolor}17}]:} float
\end{Verbatim}
            
    \begin{Verbatim}[commandchars=\\\{\}]
{\color{incolor}In [{\color{incolor}18}]:} \PY{c+c1}{\PYZsh{} le type complexe s\PYZsq{}appelle \PYZsq{}complex\PYZsq{}}
         \PY{n+nb}{type}\PY{p}{(}\PY{l+m+mi}{1}\PY{n}{j}\PY{p}{)}
\end{Verbatim}


\begin{Verbatim}[commandchars=\\\{\}]
{\color{outcolor}Out[{\color{outcolor}18}]:} complex
\end{Verbatim}
            
    \hypertarget{chauxeenes-de-caractuxe8res}{%
\subsubsection{Chaînes de
caractères}\label{chauxeenes-de-caractuxe8res}}

    On a également rapidement besoin de chaînes de caractères, on les
étudiera bientôt en détail, mais en guise d'avant-goût~:

    \begin{Verbatim}[commandchars=\\\{\}]
{\color{incolor}In [{\color{incolor}19}]:} \PY{n}{chaine} \PY{o}{=} \PY{l+s+s2}{\PYZdq{}}\PY{l+s+s2}{Bonjour le monde !}\PY{l+s+s2}{\PYZdq{}}
         \PY{n+nb}{print}\PY{p}{(}\PY{n}{chaine}\PY{p}{)}
\end{Verbatim}


    \begin{Verbatim}[commandchars=\\\{\}]
Bonjour le monde !

    \end{Verbatim}

    \hypertarget{conversions}{%
\subsubsection{Conversions}\label{conversions}}

    Il est parfois nécessaire de convertir une donnée d'un type dans un
autre. Par exemple on peut demander à l'utilisateur d'entrer une valeur
au clavier grâce à la fonction \texttt{input}, comme ceci~:

    \begin{Verbatim}[commandchars=\\\{\}]
{\color{incolor}In [{\color{incolor}20}]:} \PY{n}{reponse} \PY{o}{=} \PY{n+nb}{input}\PY{p}{(}\PY{l+s+s2}{\PYZdq{}}\PY{l+s+s2}{quel est votre âge ? }\PY{l+s+s2}{\PYZdq{}}\PY{p}{)}
         
         \PY{c+c1}{\PYZsh{} NOTE:}
         \PY{c+c1}{\PYZsh{} auto\PYZhy{}exec\PYZhy{}for\PYZhy{}latex has used instead:}
         \PY{c+c1}{\PYZsh{}\PYZsh{}\PYZsh{}\PYZsh{}\PYZsh{}\PYZsh{}\PYZsh{}\PYZsh{}\PYZsh{}\PYZsh{}}
         \PY{n}{reponse} \PY{o}{=} \PY{l+s+s1}{\PYZsq{}}\PY{l+s+s1}{25}\PY{l+s+s1}{\PYZsq{}}
         \PY{c+c1}{\PYZsh{}\PYZsh{}\PYZsh{}\PYZsh{}\PYZsh{}\PYZsh{}\PYZsh{}\PYZsh{}\PYZsh{}\PYZsh{}}
\end{Verbatim}


    \begin{Verbatim}[commandchars=\\\{\}]
{\color{incolor}In [{\color{incolor}21}]:} \PY{c+c1}{\PYZsh{} vous avez entré la chaîne suivante}
         \PY{n+nb}{print}\PY{p}{(}\PY{n}{reponse}\PY{p}{)}
\end{Verbatim}


    \begin{Verbatim}[commandchars=\\\{\}]
25

    \end{Verbatim}

    \begin{Verbatim}[commandchars=\\\{\}]
{\color{incolor}In [{\color{incolor}22}]:} \PY{c+c1}{\PYZsh{} ici reponse est une variable, et son contenu est de type chaîne de caractères}
         \PY{n+nb}{type}\PY{p}{(}\PY{n}{reponse}\PY{p}{)}
\end{Verbatim}


\begin{Verbatim}[commandchars=\\\{\}]
{\color{outcolor}Out[{\color{outcolor}22}]:} str
\end{Verbatim}
            
    Maintenant je veux faire des calculs sur votre âge, par exemple le
multiplier par 2. Si je m'y prends naïvement, ça donne ceci~:

    \begin{Verbatim}[commandchars=\\\{\}]
{\color{incolor}In [{\color{incolor}23}]:} \PY{c+c1}{\PYZsh{} multiplier une chaîne de caractères par deux ne fait pas ce que l\PYZsq{}on veut,}
         \PY{c+c1}{\PYZsh{} nous verrons plus tard que ça fait une concaténation}
         \PY{l+m+mi}{2} \PY{o}{*} \PY{n}{reponse}
\end{Verbatim}


\begin{Verbatim}[commandchars=\\\{\}]
{\color{outcolor}Out[{\color{outcolor}23}]:} '2525'
\end{Verbatim}
            
    C'est pourquoi il me faut ici d'abord \textbf{convertir} la (valeur de
la) variable \texttt{reponse} en un entier, que je peux ensuite doubler
(assurez-vous d'avoir bien entré ci-dessus une valeur qui correspond à
un nombre entier)

    \begin{Verbatim}[commandchars=\\\{\}]
{\color{incolor}In [{\color{incolor}24}]:} \PY{c+c1}{\PYZsh{} reponse est une chaine}
         \PY{c+c1}{\PYZsh{} je la convertis en entier en appelant la fonction int()}
         \PY{n}{age} \PY{o}{=} \PY{n+nb}{int}\PY{p}{(}\PY{n}{reponse}\PY{p}{)}
         \PY{n+nb}{type}\PY{p}{(}\PY{n}{age}\PY{p}{)}
\end{Verbatim}


\begin{Verbatim}[commandchars=\\\{\}]
{\color{outcolor}Out[{\color{outcolor}24}]:} int
\end{Verbatim}
            
    \begin{Verbatim}[commandchars=\\\{\}]
{\color{incolor}In [{\color{incolor}25}]:} \PY{c+c1}{\PYZsh{} que je peux maintenant multiplier par 2}
         \PY{l+m+mi}{2} \PY{o}{*} \PY{n}{age}
\end{Verbatim}


\begin{Verbatim}[commandchars=\\\{\}]
{\color{outcolor}Out[{\color{outcolor}25}]:} 50
\end{Verbatim}
            
    Ou si on préfère, en une seule fois~:

    \begin{Verbatim}[commandchars=\\\{\}]
{\color{incolor}In [{\color{incolor}26}]:} \PY{n+nb}{print}\PY{p}{(}\PY{l+s+s2}{\PYZdq{}}\PY{l+s+s2}{le double de votre age est}\PY{l+s+s2}{\PYZdq{}}\PY{p}{,} \PY{l+m+mi}{2}\PY{o}{*}\PY{n+nb}{int}\PY{p}{(}\PY{n}{reponse}\PY{p}{)}\PY{p}{)}
\end{Verbatim}


    \begin{Verbatim}[commandchars=\\\{\}]
le double de votre age est 50

    \end{Verbatim}

    \hypertarget{conversions---suite}{%
\subsubsection{Conversions - suite}\label{conversions---suite}}

    De manière plus générale, pour convertir un objet en un entier, un
flottant, ou une chaîne de caractères, on peut simplement appeler une
fonction \emph{built-in} qui porte le même nom que le type cible~:

\begin{longtable}[]{@{}ll@{}}
\toprule
Type & Fonction\tabularnewline
\midrule
\endhead
Entier & \texttt{int}\tabularnewline
Flottant & \texttt{float}\tabularnewline
Complexe & \texttt{complex}\tabularnewline
Chaîne & \texttt{str}\tabularnewline
\bottomrule
\end{longtable}

Ainsi dans l'exemple précédent, \texttt{int(reponse)} représente la
conversion de \texttt{reponse} en entier.

On a illustré cette même technique dans les exemples suivants~:

    \begin{Verbatim}[commandchars=\\\{\}]
{\color{incolor}In [{\color{incolor}27}]:} \PY{c+c1}{\PYZsh{} dans l\PYZsq{}autre sens, si j\PYZsq{}ai un entier}
         \PY{n}{a} \PY{o}{=} \PY{l+m+mi}{2345}
\end{Verbatim}


    \begin{Verbatim}[commandchars=\\\{\}]
{\color{incolor}In [{\color{incolor}28}]:} \PY{c+c1}{\PYZsh{} je peux facilement le traduire en chaîne de caractères}
         \PY{n+nb}{str}\PY{p}{(}\PY{l+m+mi}{2345}\PY{p}{)}
\end{Verbatim}


\begin{Verbatim}[commandchars=\\\{\}]
{\color{outcolor}Out[{\color{outcolor}28}]:} '2345'
\end{Verbatim}
            
    \begin{Verbatim}[commandchars=\\\{\}]
{\color{incolor}In [{\color{incolor}29}]:} \PY{c+c1}{\PYZsh{} ou en complexe}
         \PY{n+nb}{complex}\PY{p}{(}\PY{l+m+mi}{2345}\PY{p}{)}
\end{Verbatim}


\begin{Verbatim}[commandchars=\\\{\}]
{\color{outcolor}Out[{\color{outcolor}29}]:} (2345+0j)
\end{Verbatim}
            
    Nous verrons plus tard que ceci se généralise à tous les types de
Python, pour convertir un objet \texttt{x} en un type \texttt{bidule},
on appelle \texttt{bidule(x)}. On y reviendra, bien entendu.

    \hypertarget{grands-nombres}{%
\subsubsection{Grands nombres}\label{grands-nombres}}

    Comme les entiers sont de précision illimitée, on peut améliorer leur
lisibilité en insérant des caractères \texttt{\_} qui sont simplement
ignorés à l'exécution.

    \begin{Verbatim}[commandchars=\\\{\}]
{\color{incolor}In [{\color{incolor}30}]:} \PY{n}{tres\PYZus{}grand\PYZus{}nombre} \PY{o}{=} \PY{l+m+mi}{23}\PY{n}{\PYZus{}456\PYZus{}789\PYZus{}012\PYZus{}345}
         
         \PY{n}{tres\PYZus{}grand\PYZus{}nombre}
\end{Verbatim}


\begin{Verbatim}[commandchars=\\\{\}]
{\color{outcolor}Out[{\color{outcolor}30}]:} 23456789012345
\end{Verbatim}
            
    \begin{Verbatim}[commandchars=\\\{\}]
{\color{incolor}In [{\color{incolor}31}]:} \PY{c+c1}{\PYZsh{} ça marche aussi avec les flottants}
         \PY{l+m+mi}{123}\PY{n}{\PYZus{}456}\PY{o}{.}\PY{l+m+mi}{789}\PY{n}{\PYZus{}012}
\end{Verbatim}


\begin{Verbatim}[commandchars=\\\{\}]
{\color{outcolor}Out[{\color{outcolor}31}]:} 123456.789012
\end{Verbatim}
            
    \hypertarget{entiers-et-bases}{%
\subsubsection{Entiers et bases}\label{entiers-et-bases}}

    Les calculettes scientifiques permettent habituellement d'entrer les
entiers dans d'autres bases que la base 10.

En Python, on peut aussi entrer un entier sous forme binaire comme
ceci~:

    \begin{Verbatim}[commandchars=\\\{\}]
{\color{incolor}In [{\color{incolor}32}]:} \PY{n}{deux\PYZus{}cents} \PY{o}{=} \PY{l+m+mb}{0b11001000}
         \PY{n+nb}{print}\PY{p}{(}\PY{n}{deux\PYZus{}cents}\PY{p}{)}
\end{Verbatim}


    \begin{Verbatim}[commandchars=\\\{\}]
200

    \end{Verbatim}

    Ou encore sous forme octale (en base 8) comme ceci~:

    \begin{Verbatim}[commandchars=\\\{\}]
{\color{incolor}In [{\color{incolor}33}]:} \PY{n}{deux\PYZus{}cents} \PY{o}{=} \PY{l+m+mo}{0o310}
         \PY{n+nb}{print}\PY{p}{(}\PY{n}{deux\PYZus{}cents}\PY{p}{)}
\end{Verbatim}


    \begin{Verbatim}[commandchars=\\\{\}]
200

    \end{Verbatim}

    Ou enfin encore en hexadécimal (base 16) comme ceci~:

    \begin{Verbatim}[commandchars=\\\{\}]
{\color{incolor}In [{\color{incolor}34}]:} \PY{n}{deux\PYZus{}cents} \PY{o}{=} \PY{l+m+mh}{0xc8}
         \PY{n+nb}{print}\PY{p}{(}\PY{n}{deux\PYZus{}cents}\PY{p}{)}
\end{Verbatim}


    \begin{Verbatim}[commandchars=\\\{\}]
200

    \end{Verbatim}

    Pour d'autres bases, on peut utiliser la fonction de conversion
\texttt{int} en lui passant un argument supplémentaire~:

    \begin{Verbatim}[commandchars=\\\{\}]
{\color{incolor}In [{\color{incolor}35}]:} \PY{n}{deux\PYZus{}cents} \PY{o}{=} \PY{n+nb}{int}\PY{p}{(}\PY{l+s+s1}{\PYZsq{}}\PY{l+s+s1}{3020}\PY{l+s+s1}{\PYZsq{}}\PY{p}{,} \PY{l+m+mi}{4}\PY{p}{)}
         \PY{n+nb}{print}\PY{p}{(}\PY{n}{deux\PYZus{}cents}\PY{p}{)}
\end{Verbatim}


    \begin{Verbatim}[commandchars=\\\{\}]
200

    \end{Verbatim}

    \hypertarget{fonctions-mathuxe9matiques}{%
\subsubsection{Fonctions
mathématiques}\label{fonctions-mathuxe9matiques}}

    Python fournit naturellement un ensemble très complet d'opérateurs
mathématiques pour les fonctions exponentielles, trigonométriques et
autres, mais leur utilisation ne nous est pas encore accessible à ce
stade et nous les verrons ultérieurement.


    % Add a bibliography block to the postdoc
    
    
    

        
    
    
    

    

    \hypertarget{affectations-et-opuxe9rations-uxe0-la}{%
\section{\texorpdfstring{Affectations et Opérations (à la
\texttt{+=})}{Affectations et Opérations (à la +=)}}\label{affectations-et-opuxe9rations-uxe0-la}}

    \hypertarget{compluxe9ment---niveau-intermuxe9diaire}{%
\subsection{Complément - niveau
intermédiaire}\label{compluxe9ment---niveau-intermuxe9diaire}}

    Il existe en Python toute une famille d'opérateurs dérivés de
l'affectation qui permettent de faire en une fois une opération et une
affectation. En voici quelques exemples.

    \hypertarget{incruxe9mentation}{%
\subsubsection{Incrémentation}\label{incruxe9mentation}}

    On peut facilement augmenter la valeur d'une variable numérique comme
ceci~:

    \begin{Verbatim}[commandchars=\\\{\},frame=single,framerule=0.3mm,rulecolor=\color{cellframecolor}]
{\color{incolor}In [{\color{incolor}1}]:} \PY{n}{entier} \PY{o}{=} \PY{l+m+mi}{10}
        
        \PY{n}{entier} \PY{o}{+}\PY{o}{=} \PY{l+m+mi}{2}
        \PY{n+nb}{print}\PY{p}{(}\PY{l+s+s1}{\PYZsq{}}\PY{l+s+s1}{entier}\PY{l+s+s1}{\PYZsq{}}\PY{p}{,} \PY{n}{entier}\PY{p}{)}
\end{Verbatim}


    \begin{Verbatim}[commandchars=\\\{\},frame=single,framerule=0.3mm,rulecolor=\color{cellframecolor}]
entier 12
\end{Verbatim}

    Comme on le devine peut-être, ceci est équivalent à~:

    \begin{Verbatim}[commandchars=\\\{\},frame=single,framerule=0.3mm,rulecolor=\color{cellframecolor}]
{\color{incolor}In [{\color{incolor}2}]:} \PY{n}{entier} \PY{o}{=} \PY{l+m+mi}{10}
        
        \PY{n}{entier} \PY{o}{=} \PY{n}{entier} \PY{o}{+} \PY{l+m+mi}{2}
        \PY{n+nb}{print}\PY{p}{(}\PY{l+s+s1}{\PYZsq{}}\PY{l+s+s1}{entier}\PY{l+s+s1}{\PYZsq{}}\PY{p}{,} \PY{n}{entier}\PY{p}{)}
\end{Verbatim}


    \begin{Verbatim}[commandchars=\\\{\},frame=single,framerule=0.3mm,rulecolor=\color{cellframecolor}]
entier 12
\end{Verbatim}

    \hypertarget{autres-opuxe9rateurs-courants}{%
\subsubsection{Autres opérateurs
courants}\label{autres-opuxe9rateurs-courants}}

    Cette forme, qui combine opération sur une variable et réaffectation du
résultat à la même variable, est disponible avec tous les opérateurs
courants~:

    \begin{Verbatim}[commandchars=\\\{\},frame=single,framerule=0.3mm,rulecolor=\color{cellframecolor}]
{\color{incolor}In [{\color{incolor}3}]:} \PY{n}{entier} \PY{o}{\PYZhy{}}\PY{o}{=} \PY{l+m+mi}{4}
        \PY{n+nb}{print}\PY{p}{(}\PY{l+s+s1}{\PYZsq{}}\PY{l+s+s1}{après décrément}\PY{l+s+s1}{\PYZsq{}}\PY{p}{,} \PY{n}{entier}\PY{p}{)}
        \PY{n}{entier} \PY{o}{*}\PY{o}{=} \PY{l+m+mi}{2}
        \PY{n+nb}{print}\PY{p}{(}\PY{l+s+s1}{\PYZsq{}}\PY{l+s+s1}{après doublement}\PY{l+s+s1}{\PYZsq{}}\PY{p}{,} \PY{n}{entier}\PY{p}{)}
        \PY{n}{entier} \PY{o}{/}\PY{o}{=} \PY{l+m+mi}{2}
        \PY{n+nb}{print}\PY{p}{(}\PY{l+s+s1}{\PYZsq{}}\PY{l+s+s1}{mis à moitié}\PY{l+s+s1}{\PYZsq{}}\PY{p}{,} \PY{n}{entier}\PY{p}{)}
\end{Verbatim}


    \begin{Verbatim}[commandchars=\\\{\},frame=single,framerule=0.3mm,rulecolor=\color{cellframecolor}]
après décrément 8
après doublement 16
mis à moitié 8.0
\end{Verbatim}

    \hypertarget{types-non-numuxe9riques}{%
\subsubsection{Types non numériques}\label{types-non-numuxe9riques}}

    En réalité cette construction est disponible sur tous les types qui
supportent l'opérateur en question. Par exemple, les listes (que nous
verrons bientôt) peuvent être additionnées entre elles~:

    \begin{Verbatim}[commandchars=\\\{\},frame=single,framerule=0.3mm,rulecolor=\color{cellframecolor}]
{\color{incolor}In [{\color{incolor}4}]:} \PY{n}{liste} \PY{o}{=} \PY{p}{[}\PY{l+m+mi}{0}\PY{p}{,} \PY{l+m+mi}{3}\PY{p}{,} \PY{l+m+mi}{5}\PY{p}{]}
        \PY{n+nb}{print}\PY{p}{(}\PY{l+s+s1}{\PYZsq{}}\PY{l+s+s1}{liste}\PY{l+s+s1}{\PYZsq{}}\PY{p}{,} \PY{n}{liste}\PY{p}{)}
        
        \PY{n}{liste} \PY{o}{+}\PY{o}{=} \PY{p}{[}\PY{l+s+s1}{\PYZsq{}}\PY{l+s+s1}{a}\PY{l+s+s1}{\PYZsq{}}\PY{p}{,} \PY{l+s+s1}{\PYZsq{}}\PY{l+s+s1}{b}\PY{l+s+s1}{\PYZsq{}}\PY{p}{]}
        \PY{n+nb}{print}\PY{p}{(}\PY{l+s+s1}{\PYZsq{}}\PY{l+s+s1}{après ajout}\PY{l+s+s1}{\PYZsq{}}\PY{p}{,} \PY{n}{liste}\PY{p}{)}
\end{Verbatim}


    \begin{Verbatim}[commandchars=\\\{\},frame=single,framerule=0.3mm,rulecolor=\color{cellframecolor}]
liste [0, 3, 5]
après ajout [0, 3, 5, 'a', 'b']
\end{Verbatim}

    Beaucoup de types supportent l'opérateur \texttt{+}, qui est sans doute
de loin celui qui est le plus utilisé avec cette construction.

    \hypertarget{opuxe9rateurs-plus-abscons}{%
\subsubsection{Opérateurs plus
abscons}\label{opuxe9rateurs-plus-abscons}}

    Signalons enfin que l'on trouve aussi cette construction avec d'autres
opérateurs moins fréquents, par exemple~:

    \begin{Verbatim}[commandchars=\\\{\},frame=single,framerule=0.3mm,rulecolor=\color{cellframecolor}]
{\color{incolor}In [{\color{incolor}5}]:} \PY{n}{entier} \PY{o}{=} \PY{l+m+mi}{2}
        \PY{n+nb}{print}\PY{p}{(}\PY{l+s+s1}{\PYZsq{}}\PY{l+s+s1}{entier:}\PY{l+s+s1}{\PYZsq{}}\PY{p}{,} \PY{n}{entier}\PY{p}{)}
        \PY{n}{entier} \PY{o}{*}\PY{o}{*}\PY{o}{=} \PY{l+m+mi}{10}
        \PY{n+nb}{print}\PY{p}{(}\PY{l+s+s1}{\PYZsq{}}\PY{l+s+s1}{à la puissance dix:}\PY{l+s+s1}{\PYZsq{}}\PY{p}{,} \PY{n}{entier}\PY{p}{)}
        \PY{n}{entier} \PY{o}{\PYZpc{}}\PY{o}{=} \PY{l+m+mi}{5}
        \PY{n+nb}{print}\PY{p}{(}\PY{l+s+s1}{\PYZsq{}}\PY{l+s+s1}{modulo 5:}\PY{l+s+s1}{\PYZsq{}}\PY{p}{,} \PY{n}{entier}\PY{p}{)}
\end{Verbatim}


    \begin{Verbatim}[commandchars=\\\{\},frame=single,framerule=0.3mm,rulecolor=\color{cellframecolor}]
entier: 2
à la puissance dix: 1024
modulo 5: 4
\end{Verbatim}

    Et pour ceux qui connaissent déjà un peu Python, on peut même le faire
avec des opérateurs de décalage, que nous verrons très bientôt~:

    \begin{Verbatim}[commandchars=\\\{\},frame=single,framerule=0.3mm,rulecolor=\color{cellframecolor}]
{\color{incolor}In [{\color{incolor}6}]:} \PY{n}{entier} \PY{o}{\PYZlt{}\PYZlt{}}\PY{o}{=} \PY{l+m+mi}{2}
        \PY{n+nb}{print}\PY{p}{(}\PY{l+s+s1}{\PYZsq{}}\PY{l+s+s1}{double décalage gauche:}\PY{l+s+s1}{\PYZsq{}}\PY{p}{,} \PY{n}{entier}\PY{p}{)}
\end{Verbatim}


    \begin{Verbatim}[commandchars=\\\{\},frame=single,framerule=0.3mm,rulecolor=\color{cellframecolor}]
double décalage gauche: 16
\end{Verbatim}


    % Add a bibliography block to the postdoc
    
    
    

        
    
    
    

    

    \hypertarget{notions-sur-la-pruxe9cision-des-calculs-flottants}{%
\section{Notions sur la précision des calculs
flottants}\label{notions-sur-la-pruxe9cision-des-calculs-flottants}}

    \hypertarget{compluxe9ment---niveau-avancuxe9}{%
\subsection{Complément - niveau
avancé}\label{compluxe9ment---niveau-avancuxe9}}

    \hypertarget{le-probluxe8me}{%
\subsubsection{Le problème}\label{le-probluxe8me}}

    Comme pour les entiers, les calculs sur les flottants sont,
naturellement, réalisés par le processeur. Cependant contrairement au
cas des entiers où les calculs sont toujours exacts, les flottants
posent un problème de précision. Cela n'est pas propre au langage
Python, mais est dû à la technique de codage des nombres flottants sous
forme binaire.

    Voyons tout d'abord comment se matérialise le problème~:

    \begin{Verbatim}[commandchars=\\\{\},frame=single,framerule=0.3mm,rulecolor=\color{cellframecolor}]
{\color{incolor}In [{\color{incolor}1}]:} \PY{l+m+mf}{0.2} \PY{o}{+} \PY{l+m+mf}{0.4}
\end{Verbatim}


\begin{Verbatim}[commandchars=\\\{\},frame=single,framerule=0.3mm,rulecolor=\color{cellframecolor}]
{\color{outcolor}Out[{\color{outcolor}1}]:} 0.6000000000000001
\end{Verbatim}
            
    Il faut retenir que lorsqu'on écrit un nombre flottant sous forme
décimale, la valeur utilisée en mémoire pour représenter ce nombre,
parce que cette valeur est codée en binaire, ne représente \textbf{pas
toujours exactement} le nombre entré.

    \begin{Verbatim}[commandchars=\\\{\},frame=single,framerule=0.3mm,rulecolor=\color{cellframecolor}]
{\color{incolor}In [{\color{incolor}2}]:} \PY{c+c1}{\PYZsh{} du coup cette expression est fausse, à cause de l\PYZsq{}erreur d\PYZsq{}arrondi}
        \PY{l+m+mf}{0.3} \PY{o}{\PYZhy{}} \PY{l+m+mf}{0.1} \PY{o}{==} \PY{l+m+mf}{0.2}
\end{Verbatim}


\begin{Verbatim}[commandchars=\\\{\},frame=single,framerule=0.3mm,rulecolor=\color{cellframecolor}]
{\color{outcolor}Out[{\color{outcolor}2}]:} False
\end{Verbatim}
            
    Aussi, comme on le voit, les différentes erreurs d'arrondi qui se
produisent à chaque étape du calcul s'accumulent et produisent un
résultat parfois surprenant. De nouveau, ce problème n'est pas
spécifique à Python, il existe pour tous les langages, et il est bien
connu des numériciens.

    Dans une grande majorité des cas, ces erreurs d'arrondi ne sont pas
pénalisantes. Il faut toutefois en être conscient car cela peut
expliquer des comportements curieux.

    \hypertarget{une-solution-penser-en-termes-de-nombres-rationnels}{%
\subsubsection{Une solution~: penser en termes de nombres
rationnels}\label{une-solution-penser-en-termes-de-nombres-rationnels}}

    Tout d'abord si votre problème se pose bien en termes de nombres
rationnels, il est alors tout à fait possible de le résoudre avec
exactitude.

    Alors qu'il n'est pas possible d'écrire exactement \(3/10\) en base 2,
ni d'ailleurs \(1/3\) en base 10, on peut représenter
\textbf{exactement} ces nombres dès lors qu'on les considère comme des
fractions et qu'on les encode avec deux nombres entiers.

    Python fournit en standard le module \texttt{fractions} qui permet de
résoudre le problème. Voici comment on pourrait l'utiliser pour
vérifier, cette fois avec succès, que \(0.3 - 0.1\) vaut bien \(0.2\).
Ce code anticipe sur l'utilisation des modules et des classes en Python,
ici nous créons des objets de type \texttt{Fraction}~:

    \begin{Verbatim}[commandchars=\\\{\},frame=single,framerule=0.3mm,rulecolor=\color{cellframecolor}]
{\color{incolor}In [{\color{incolor}3}]:} \PY{c+c1}{\PYZsh{} on importe le module fractions, qui lui\PYZhy{}même définit le symbole Fraction}
        \PY{k+kn}{from} \PY{n+nn}{fractions} \PY{k}{import} \PY{n}{Fraction}
        
        \PY{c+c1}{\PYZsh{} et cette fois, les calculs sont exacts, et l\PYZsq{}expression retourne bien True}
        \PY{n}{Fraction}\PY{p}{(}\PY{l+m+mi}{3}\PY{p}{,} \PY{l+m+mi}{10}\PY{p}{)} \PY{o}{\PYZhy{}} \PY{n}{Fraction}\PY{p}{(}\PY{l+m+mi}{1}\PY{p}{,} \PY{l+m+mi}{10}\PY{p}{)} \PY{o}{==} \PY{n}{Fraction}\PY{p}{(}\PY{l+m+mi}{2}\PY{p}{,} \PY{l+m+mi}{10}\PY{p}{)}
\end{Verbatim}


\begin{Verbatim}[commandchars=\\\{\},frame=single,framerule=0.3mm,rulecolor=\color{cellframecolor}]
{\color{outcolor}Out[{\color{outcolor}3}]:} True
\end{Verbatim}
            
    Ou encore d'ailleurs, équivalent et plus lisible~:

    \begin{Verbatim}[commandchars=\\\{\},frame=single,framerule=0.3mm,rulecolor=\color{cellframecolor}]
{\color{incolor}In [{\color{incolor}4}]:} \PY{n}{Fraction}\PY{p}{(}\PY{l+s+s1}{\PYZsq{}}\PY{l+s+s1}{0.3}\PY{l+s+s1}{\PYZsq{}}\PY{p}{)} \PY{o}{\PYZhy{}} \PY{n}{Fraction}\PY{p}{(}\PY{l+s+s1}{\PYZsq{}}\PY{l+s+s1}{0.1}\PY{l+s+s1}{\PYZsq{}}\PY{p}{)} \PY{o}{==} \PY{n}{Fraction}\PY{p}{(}\PY{l+s+s1}{\PYZsq{}}\PY{l+s+s1}{2/10}\PY{l+s+s1}{\PYZsq{}}\PY{p}{)}
\end{Verbatim}


\begin{Verbatim}[commandchars=\\\{\},frame=single,framerule=0.3mm,rulecolor=\color{cellframecolor}]
{\color{outcolor}Out[{\color{outcolor}4}]:} True
\end{Verbatim}
            
    \hypertarget{une-autre-solution-le-module-decimal}{%
\subsubsection{Une autre solution~: le module
decimal}\label{une-autre-solution-le-module-decimal}}

    Si par contre vous ne manipulez pas des nombres rationnels et que du
coup la représentation sous forme de fractions ne peut pas convenir dans
votre cas, signalons l'existence du module standard \texttt{decimal} qui
offre des fonctionnalités très voisines du type \texttt{float}, tout en
éliminant la plupart des inconvénients, au prix naturellement d'une
consommation mémoire supérieure.

    Pour reprendre l'exemple de départ, mais en utilisant le module decimal,
on écrirait alors~:

    \begin{Verbatim}[commandchars=\\\{\},frame=single,framerule=0.3mm,rulecolor=\color{cellframecolor}]
{\color{incolor}In [{\color{incolor}5}]:} \PY{k+kn}{from} \PY{n+nn}{decimal} \PY{k}{import} \PY{n}{Decimal}
        
        \PY{n}{Decimal}\PY{p}{(}\PY{l+s+s1}{\PYZsq{}}\PY{l+s+s1}{0.3}\PY{l+s+s1}{\PYZsq{}}\PY{p}{)} \PY{o}{\PYZhy{}} \PY{n}{Decimal}\PY{p}{(}\PY{l+s+s1}{\PYZsq{}}\PY{l+s+s1}{0.1}\PY{l+s+s1}{\PYZsq{}}\PY{p}{)} \PY{o}{==} \PY{n}{Decimal}\PY{p}{(}\PY{l+s+s1}{\PYZsq{}}\PY{l+s+s1}{0.2}\PY{l+s+s1}{\PYZsq{}}\PY{p}{)}
\end{Verbatim}


\begin{Verbatim}[commandchars=\\\{\},frame=single,framerule=0.3mm,rulecolor=\color{cellframecolor}]
{\color{outcolor}Out[{\color{outcolor}5}]:} True
\end{Verbatim}
            
    \hypertarget{pour-aller-plus-loin}{%
\subsubsection{Pour aller plus loin}\label{pour-aller-plus-loin}}

    Tous ces documents sont en anglais~:

\begin{itemize}
\tightlist
\item
  un
  \href{https://docs.python.org/3/tutorial/floatingpoint.html}{tutoriel
  sur les nombres flottants}~;
\item
  la
  \href{https://docs.python.org/3/library/fractions.html}{documentation
  sur la classe Fraction}~;
\item
  la \href{https://docs.python.org/3/library/decimal.html}{documentation
  sur la classe Decimal}.
\end{itemize}


    % Add a bibliography block to the postdoc
    
    
    

        \hypertarget{opuxe9rations-bit-uxe0-bit-bitwise}{%
\section{\texorpdfstring{Opérations \emph{bit à bit}
(\emph{bitwise})}{Opérations bit à bit (bitwise)}}\label{opuxe9rations-bit-uxe0-bit-bitwise}}

    \hypertarget{compluxe9ments---niveau-avancuxe9}{%
\subsection{Compléments - niveau
avancé}\label{compluxe9ments---niveau-avancuxe9}}

    Les compléments ci-dessous expliquent des fonctions évoluées sur les
entiers. Les débutants en programmation peuvent sans souci sauter cette
partie en cas de difficultés.

    \hypertarget{opuxe9rations-logiques-et-ou-et-ou-exclusif}{%
\subsubsection{\texorpdfstring{Opérations logiques~: \emph{ET}
\texttt{\&}, \emph{OU} \texttt{\textbar{}} et \emph{OU} exclusif
\texttt{\^{}}}{Opérations logiques~: ET \&, OU \textbar{} et OU exclusif \^{}}}\label{opuxe9rations-logiques-et-ou-et-ou-exclusif}}

    Il est possible aussi de faire des opérations \emph{bit à bit} sur les
nombres entiers. Le plus simple est de penser à l'écriture du nombre en
base 2.\\

Considérons par exemple deux entiers constants dans cet exercice

    \begin{Verbatim}[commandchars=\\\{\}]
{\color{incolor}In [{\color{incolor}1}]:} \PY{n}{x49} \PY{o}{=} \PY{l+m+mi}{49}
        \PY{n}{y81} \PY{o}{=} \PY{l+m+mi}{81}
\end{Verbatim}


    Ce qui nous donne comme décomposition binaire~:\\

\(\begin{array}{rcccccc} x49 & = & 49 & = & 32 + 16 + 1 & \rightarrow &(0,1,1,0,0,0,1) \\ y81 & = & 81 & = & 64 + 16 + 1 & \rightarrow &(1,0,1,0,0,0,1) \end{array}\)\\

    Pour comprendre comment passer de \(32 + 16 + 1\) à \((0,1,1,0,0,0,1)\)
il suffit d'observer que~:\\

\(32 + 16 + 1 = \textbf{0}*2^6 + \textbf{1}*2^5 + \textbf{1}*2^4 + \textbf{0}*2^3 + \textbf{0}*2^2 + \textbf{0}*2^1 + \textbf{1}*2^0\)

    \hypertarget{et-logique-opuxe9rateur}{%
\paragraph{\texorpdfstring{\emph{ET} logique~: opérateur
\texttt{\&}}{ET logique~: opérateur \&}\\\\}\label{et-logique-opuxe9rateur}}

    L'opération logique \texttt{\&} va faire un \emph{et} logique bit à bit
entre les opérandes, ainsi

    \begin{Verbatim}[commandchars=\\\{\}]
{\color{incolor}In [{\color{incolor}2}]:} \PY{n}{x49} \PY{o}{\PYZam{}} \PY{n}{y81}
\end{Verbatim}


\begin{Verbatim}[commandchars=\\\{\}]
{\color{outcolor}Out[{\color{outcolor}2}]:} 17
\end{Verbatim}
            
    Et en effet~:\\

\(\begin{array}{rcl} x49 & \rightarrow & (0,1,1,0,0,0,1) \\ y81 & \rightarrow & (1,0,1,0,0,0,1) \\ x49\ \&\ y81 & \rightarrow & (0,0,1,0,0,0,1) \rightarrow 16 + 1 \rightarrow 17 \end{array}\)

    \hypertarget{ou-logique-opuxe9rateur}{%
\paragraph{\texorpdfstring{\emph{OU} logique~: opérateur
\texttt{\textbar{}}}{OU logique~: opérateur \textbar{}}\\\\}\label{ou-logique-opuxe9rateur}}

    De même, l'opérateur logique \texttt{\textbar{}} fait simplement un
\emph{ou} logique, comme ceci~:

    \begin{Verbatim}[commandchars=\\\{\}]
{\color{incolor}In [{\color{incolor}3}]:} \PY{n}{x49} \PY{o}{|} \PY{n}{y81}
\end{Verbatim}


\begin{Verbatim}[commandchars=\\\{\}]
{\color{outcolor}Out[{\color{outcolor}3}]:} 113
\end{Verbatim}
            
    On s'y retrouve parce que~:\\

\(\begin{array}{rcl} x49 & \rightarrow & (0,1,1,0,0,0,1) \\ y81 & \rightarrow & (1,0,1,0,0,0,1) \\ x49\ |\ y81 & \rightarrow & (1,1,1,0,0,0,1) \rightarrow 64 + 32 + 16 + 1 \rightarrow 113 \end{array}\)

    \hypertarget{ou-exclusif-opuxe9rateur}{%
\paragraph{\texorpdfstring{\emph{OU} exclusif~: opérateur
\texttt{\^{}}}{OU exclusif~: opérateur \^{}}\\\\}\label{ou-exclusif-opuxe9rateur}}

    Enfin, on peut également faire la même opération à base de \emph{ou}
exclusif avec l'opérateur \texttt{\^{}}~:

    \begin{Verbatim}[commandchars=\\\{\}]
{\color{incolor}In [{\color{incolor}4}]:} \PY{n}{x49} \PY{o}{\PYZca{}} \PY{n}{y81}
\end{Verbatim}


\begin{Verbatim}[commandchars=\\\{\}]
{\color{outcolor}Out[{\color{outcolor}4}]:} 96
\end{Verbatim}
            
    Je vous laisse le soin de décortiquer le calcul à titre d'exercice (le
\emph{ou} exclusif de deux bits est vrai si et seulement si exactement
une des deux entrées est vraie).

    \hypertarget{duxe9calages}{%
\subsubsection{Décalages}\label{duxe9calages}}

    Un décalage \textbf{à gauche} de, par exemple, 4 positions, revient à
décaler tout le champ de bits de 4 cases à gauche (les 4 nouveaux bits
insérés sont toujours des 0). C'est donc équivalent à une
\textbf{multiplication} par \(2^4 = 16\)~:

    \begin{Verbatim}[commandchars=\\\{\}]
{\color{incolor}In [{\color{incolor}5}]:} \PY{n}{x49} \PY{o}{\PYZlt{}\PYZlt{}} \PY{l+m+mi}{4}
\end{Verbatim}


\begin{Verbatim}[commandchars=\\\{\}]
{\color{outcolor}Out[{\color{outcolor}5}]:} 784
\end{Verbatim}
            
    \(\begin{array}{rcl} x49 & \rightarrow & (0,1,1,0,0,0,1) \\ x49\ <<\ 4 & \rightarrow & (0,1,1,0,0,0,1,0,0,0,0) \rightarrow 512 + 256 + 16 \rightarrow 784 \end{array}\)

    De la même manière, le décalage à \textbf{droite} de \(n\) revient à une
\textbf{division} par \(2^n\) (plus précisément, le quotient de la
division)~:

    \begin{Verbatim}[commandchars=\\\{\}]
{\color{incolor}In [{\color{incolor}6}]:} \PY{n}{x49} \PY{o}{\PYZgt{}\PYZgt{}} \PY{l+m+mi}{4}
\end{Verbatim}


\begin{Verbatim}[commandchars=\\\{\}]
{\color{outcolor}Out[{\color{outcolor}6}]:} 3
\end{Verbatim}
            
    \(\begin{array}{rcl} x49 & \rightarrow & (0,1,1,0,0,0,1) \\ x49\ >>\ 4 & \rightarrow & (0,0,0,0,0,1,1) \rightarrow 2 + 1 \rightarrow 3 \end{array}\)

    \hypertarget{une-astuce}{%
\subsubsection{Une astuce}\label{une-astuce}}

    On peut utiliser la fonction \emph{built-in} \texttt{bin} pour calculer
la représentation binaire d'un entier. Attention, la valeur de retour
est une chaîne de caractères de type \texttt{str}~:

    \begin{Verbatim}[commandchars=\\\{\}]
{\color{incolor}In [{\color{incolor}7}]:} \PY{n+nb}{bin}\PY{p}{(}\PY{n}{x49}\PY{p}{)}
\end{Verbatim}


\begin{Verbatim}[commandchars=\\\{\}]
{\color{outcolor}Out[{\color{outcolor}7}]:} '0b110001'
\end{Verbatim}
            
    Dans l'autre sens, on peut aussi entrer un entier directement en base 2
comme ceci~:

    \begin{Verbatim}[commandchars=\\\{\}]
{\color{incolor}In [{\color{incolor}8}]:} \PY{n}{x49bis} \PY{o}{=} \PY{l+m+mb}{0b110001}
        \PY{n}{x49bis} \PY{o}{==} \PY{n}{x49}
\end{Verbatim}


\begin{Verbatim}[commandchars=\\\{\}]
{\color{outcolor}Out[{\color{outcolor}8}]:} True
\end{Verbatim}
            
    Ici, comme on le voit, \texttt{x49bis} est bien un entier.

    \hypertarget{pour-en-savoir-plus}{%
\subsubsection{Pour en savoir plus}\label{pour-en-savoir-plus}}

    \href{https://docs.python.org/3/library/stdtypes.html\#bitwise-operations-on-integer-types}{Section
de la documentation Python}.
        \hypertarget{estimer-le-plus-petit-grand-flottant}{%
\section{Estimer le plus petit (grand)
flottant}\label{estimer-le-plus-petit-grand-flottant}}

    \hypertarget{exercice---niveau-basique}{%
\subsection{Exercice - niveau basique}\label{exercice---niveau-basique}}

    \hypertarget{le-plus-petit-flottant}{%
\subsubsection{Le plus petit flottant}\label{le-plus-petit-flottant}}

    En corollaire de la discussion sur la précision des flottants, il faut
savoir que le système de codage en mémoire impose aussi une limite. Les
réels très petits, ou très grands, ne peuvent plus être représentés de
cette manière.\\

C'est notamment très gênant si vous implémentez un logiciel
probabiliste, comme des graphes de Markov, où les probabilités
d'occurrence de séquences très longues tendent très rapidement vers des
valeurs extrêmement petites.\\

    Le but de cet exercice est d'estimer la valeur du plus petit flottant
qui peut être représenté comme un flottant. Pour vous aider, voici deux
valeurs~:

    \begin{Verbatim}[commandchars=\\\{\}]
{\color{incolor}In [{\color{incolor}1}]:} \PY{l+m+mi}{10}\PY{o}{*}\PY{o}{*}\PY{o}{\PYZhy{}}\PY{l+m+mi}{320}
\end{Verbatim}


\begin{Verbatim}[commandchars=\\\{\}]
{\color{outcolor}Out[{\color{outcolor}1}]:} 1e-320
\end{Verbatim}
            
    \begin{Verbatim}[commandchars=\\\{\}]
{\color{incolor}In [{\color{incolor}2}]:} \PY{l+m+mi}{10}\PY{o}{*}\PY{o}{*}\PY{o}{\PYZhy{}}\PY{l+m+mi}{330}
\end{Verbatim}


\begin{Verbatim}[commandchars=\\\{\}]
{\color{outcolor}Out[{\color{outcolor}2}]:} 0.0
\end{Verbatim}
            
    Comme on le voit, \(10^{-320}\) est correctement imprimé, alors que
\(10^{-330}\) est, de manière erronée, rapporté comme étant nul.

    \textbf{Notes~:}

\begin{itemize}
\item
  À ce stade du cours, pour estimer le plus petit flottant, procédez
  simplement par approximations successives.
\item
  Sans utiliser de boucle, la précision que vous pourrez obtenir n'est
  que fonction de votre patience, ne dépassez pas 4 à 5 itérations
  successives :)
\item
  Il est par contre pertinent d'utiliser une approche rationnelle pour
  déterminer l'itération suivante (par opposition à une approche ``au
  petit bonheur''). Pour ceux qui ne connaissent pas, nous vous
  recommandons de vous documenter sur l'algorithme de
  \href{https://fr.wikipedia.org/wiki/Recherche_dichotomique}{\textbf{dichotomie}}.
\end{itemize}

    \begin{Verbatim}[commandchars=\\\{\}]
{\color{incolor}In [{\color{incolor}3}]:} \PY{l+m+mi}{10}\PY{o}{*}\PY{o}{*}\PY{o}{\PYZhy{}}\PY{l+m+mi}{325}
\end{Verbatim}


\begin{Verbatim}[commandchars=\\\{\}]
{\color{outcolor}Out[{\color{outcolor}3}]:} 0.0
\end{Verbatim}
            
    Voici quelques cellules de code vides~; vous pouvez en créer d'autres si
nécessaire, le plus simple étant de taper \texttt{Alt+Enter}, ou
d'utiliser le menu \emph{``Insert~-\textgreater{}~Insert~Cell~Below''}

    \begin{Verbatim}[commandchars=\\\{\}]
{\color{incolor}In [{\color{incolor} }]:} \PY{c+c1}{\PYZsh{} vos essais successifs ici}
\end{Verbatim}


    \begin{Verbatim}[commandchars=\\\{\}]
{\color{incolor}In [{\color{incolor}10}]:} \PY{o}{.}\PY{l+m+mi}{24}\PY{o}{*}\PY{l+m+mi}{10}\PY{o}{*}\PY{o}{*}\PY{o}{\PYZhy{}}\PY{l+m+mi}{323}
\end{Verbatim}


\begin{Verbatim}[commandchars=\\\{\}]
{\color{outcolor}Out[{\color{outcolor}10}]:} 0.0
\end{Verbatim}
            
    \hypertarget{le-plus-grand-flottant}{%
\subsubsection{Le plus grand flottant}\label{le-plus-grand-flottant}}

    La même limitation s'applique sur les grands nombres. Toutefois, cela
est un peu moins évident, car comme toujours il faut faire attention aux
types~:

    \begin{Verbatim}[commandchars=\\\{\}]
{\color{incolor}In [{\color{incolor}4}]:} \PY{l+m+mi}{10}\PY{o}{*}\PY{o}{*}\PY{l+m+mi}{450}
\end{Verbatim}


\begin{Verbatim}[commandchars=\\\{\}]
{\color{outcolor}Out[{\color{outcolor}4}]:} 100000000000000000000000000000000000000000000000000000000000000000000000\\
	000000000000000000000000000000000000000000000000000000000000000000000000\\
	000000000000000000000000000000000000000000000000000000000000000000000000\\
	000000000000000000000000000000000000000000000000000000000000000000000000\\
	000000000000000000000000000000000000000000000000000000000000000000000000\\
	000000000000000000000000000000000000000000000000000000000000000000000000\\
	0000000000000000000
\end{Verbatim}
            
    Ce qui passe très bien car j'ai utilisé un \texttt{int} pour l'exposant,
dans ce premier cas Python calcule le résultat comme un \texttt{int},
qui est un type qui n'a pas de limitation de précision (Python utilise
intelligemment autant de bits que nécessaire pour ce genre de calculs).\\

Par contre, si j'essaie de faire le même calcul avec un exposant
flottant, Python essaie cette fois de faire son calcul avec un flottant,
et là on obtient une erreur~:

    \begin{Verbatim}[commandchars=\\\{\}]
{\color{incolor}In [{\color{incolor}5}]:} \PY{l+m+mi}{10}\PY{o}{*}\PY{o}{*}\PY{l+m+mf}{450.0}
\end{Verbatim}


    \begin{Verbatim}[commandchars=\\\{\}]

        ---------------------------------------------------------------------------

        OverflowError                             Traceback (most recent call last)

        <ipython-input-5-d63f7e475389> in <module>()
    ----> 1 10**450.0
    

        OverflowError: (34, 'Numerical result out of range')

    \end{Verbatim}

    On peut d'ailleurs remarquer que le comportement ici n'est pas
extrêmement cohérent, car avec les petits nombres Python nous a
silencieusement transformé \(10^{-330}\) en \(0\), alors que pour les
grands nombres, il lève une exception (nous verrons les exceptions plus
tard, mais vous pouvez dès maintenant remarquer que le comportement est
différent dans les deux cas).\\

    Quoi qu'il en soit, la limite pour les grands nombres se situe entre les
deux valeurs \(10^{300}\) et \(10^{310}\). On vous demande à nouveau
d'estimer comme ci-dessus une valeur approchée du plus grand nombre
qu'il soit possible de représenter comme un flottant.

    \begin{Verbatim}[commandchars=\\\{\}]
{\color{incolor}In [{\color{incolor}6}]:} \PY{l+m+mi}{10}\PY{o}{*}\PY{o}{*}\PY{l+m+mf}{300.}
\end{Verbatim}


\begin{Verbatim}[commandchars=\\\{\}]
{\color{outcolor}Out[{\color{outcolor}6}]:} 1e+300
\end{Verbatim}
            
    \begin{Verbatim}[commandchars=\\\{\}]
{\color{incolor}In [{\color{incolor}7}]:} \PY{l+m+mi}{10}\PY{o}{*}\PY{o}{*}\PY{l+m+mf}{310.}
\end{Verbatim}


    \begin{Verbatim}[commandchars=\\\{\}]

        ---------------------------------------------------------------------------

        OverflowError                             Traceback (most recent call last)

        <ipython-input-7-5d701e1fa38c> in <module>()
    ----> 1 10**310.
    

        OverflowError: (34, 'Numerical result out of range')

    \end{Verbatim}

    \begin{Verbatim}[commandchars=\\\{\}]
{\color{incolor}In [{\color{incolor} }]:} \PY{c+c1}{\PYZsh{} vos essais successifs ici}
\end{Verbatim}


    \hypertarget{compluxe9ment---niveau-avancuxe9}{%
\subsection{Complément - niveau
avancé}\label{compluxe9ment---niveau-avancuxe9}}

    En fait, on peut accéder à ces valeurs minimales et maximales pour les
flottants comme ceci

    \begin{Verbatim}[commandchars=\\\{\}]
{\color{incolor}In [{\color{incolor}8}]:} \PY{k+kn}{import} \PY{n+nn}{sys}
        \PY{n+nb}{print}\PY{p}{(}\PY{n}{sys}\PY{o}{.}\PY{n}{float\PYZus{}info}\PY{p}{)}
\end{Verbatim}


    \begin{Verbatim}[commandchars=\\\{\}]
sys.float\_info(max=1.7976931348623157e+308, max\_exp=1024, max\_10\_exp=308,
min=2.2250738585072014e-308, min\_exp=-1021, min\_10\_exp=-307, dig=15, mant\_dig=53,
epsilon=2.220446049250313e-16, radix=2, rounds=1)

    \end{Verbatim}

    Et notamment,
\href{https://docs.python.org/3/library/sys.html\#sys.float_info}{comme
expliqué ici}.

    \begin{Verbatim}[commandchars=\\\{\}]
{\color{incolor}In [{\color{incolor}9}]:} \PY{n+nb}{print}\PY{p}{(}\PY{l+s+s2}{\PYZdq{}}\PY{l+s+s2}{Flottant minimum}\PY{l+s+s2}{\PYZdq{}}\PY{p}{,} \PY{n}{sys}\PY{o}{.}\PY{n}{float\PYZus{}info}\PY{o}{.}\PY{n}{min}\PY{p}{)}
        \PY{n+nb}{print}\PY{p}{(}\PY{l+s+s2}{\PYZdq{}}\PY{l+s+s2}{Flottant maximum}\PY{l+s+s2}{\PYZdq{}}\PY{p}{,} \PY{n}{sys}\PY{o}{.}\PY{n}{float\PYZus{}info}\PY{o}{.}\PY{n}{max}\PY{p}{)}
\end{Verbatim}


    \begin{Verbatim}[commandchars=\\\{\}]
Flottant minimum 2.2250738585072014e-308
Flottant maximum 1.7976931348623157e+308

    \end{Verbatim}

    \textbf{Sauf que} vous devez avoir trouvé un maximum voisin de cette
valeur, mais le minimum observé expérimentalement ne correspond pas bien
à cette valeur.\\

Pour ceux que cela intéresse, l'explication à cette apparente
contradiction réside dans l'utilisation de
\href{http://en.wikipedia.org/wiki/Denormal\%5Fnumber}{nombres
dénormaux}.


    % Add a bibliography block to the postdoc
    
    
    
    \end{document}
